\documentclass[11pt]{article}
\usepackage[T1]{fontenc}
\usepackage{lmodern}
\usepackage[margin=1.05in]{geometry}
\usepackage{amsmath,amssymb,amsthm,mathtools}
\usepackage{microtype,enumitem,booktabs,aliascnt,flafter}
\usepackage[dvipsnames]{xcolor}
\usepackage{tikz}
\usetikzlibrary{arrows.meta,calc,positioning}
\usepackage{needspace}
\usepackage{hyperref}
\usepackage[nameinlink,noabbrev,capitalise]{cleveref}
\hypersetup{colorlinks=true,linkcolor=MidnightBlue,citecolor=MidnightBlue,
 urlcolor=MidnightBlue,pdftitle={A Modulus Proof of Cannon’s Conjecture},
 pdfauthor={OpenAI},bookmarksdepth=2}
\numberwithin{equation}{section}
\newtheorem{theorem}{Theorem}[section]
\newaliascnt{proposition}{theorem}
\newtheorem{proposition}[proposition]{Proposition}
\aliascntresetthe{proposition}
\newaliascnt{lemma}{theorem}
\newtheorem{lemma}[lemma]{Lemma}
\aliascntresetthe{lemma}
\newaliascnt{corollary}{theorem}
\newtheorem{corollary}[corollary]{Corollary}
\aliascntresetthe{corollary}
\theoremstyle{definition}
\newaliascnt{definition}{theorem}
\newtheorem{definition}[definition]{Definition}
\aliascntresetthe{definition}
\theoremstyle{remark}
\newaliascnt{remark}{theorem}
\newtheorem{remark}[remark]{Remark}
\aliascntresetthe{remark}
\crefname{lemma}{Lemma}{Lemmas}
\Crefname{lemma}{Lemma}{Lemmas}
\crefname{proposition}{Proposition}{Propositions}
\Crefname{proposition}{Proposition}{Propositions}
\crefname{corollary}{Corollary}{Corollaries}
\Crefname{corollary}{Corollary}{Corollaries}
\crefname{definition}{Definition}{Definitions}
\Crefname{definition}{Definition}{Definitions}
\crefname{remark}{Remark}{Remarks}
\Crefname{remark}{Remark}{Remarks}
\newcommand{\RR}{\mathbb R}
\newcommand{\NN}{\mathbb N}
\newcommand{\ZZ}{\mathbb Z}
\newcommand{\HH}{\mathbb H}
\newcommand{\Sphere}{\mathbb S^2}
\newcommand{\cF}{\mathcal F}
\newcommand{\cQ}{\mathcal Q}
\newcommand{\Qhat}{\widehat{\mathcal Q}}
\DeclareMathOperator{\Mod}{Mod}
\DeclareMathOperator{\osc}{osc}
\DeclareMathOperator{\diam}{diam}
\DeclareMathOperator{\dist}{dist}
\DeclareMathOperator{\id}{id}
\DeclareMathOperator{\Isom}{Isom}
\DeclareMathOperator{\Stab}{Stab}
\setlist{topsep=4pt,itemsep=2pt,parsep=0pt}
\setlength{\emergencystretch}{1em}

\title{A Modulus Proof of Cannon's Conjecture}
\author{OpenAI}
\date{September 23, 2026}

\begin{document}
\maketitle
\begin{abstract}
We prove that every hyperbolic group whose boundary is homeomorphic to
the two-sphere admits a proper cocompact isometric action on hyperbolic
three-space with finite kernel. This resolves Cannon's conjecture
positively.
\end{abstract}

\setcounter{tocdepth}{1}
\tableofcontents
\newpage

\section{Introduction}
\label{sec:introduction}

Cannon's conjecture asks whether a hyperbolic group whose boundary is a
topological two-sphere acts geometrically on hyperbolic three-space.
It links the topology at infinity of a group to its three-dimensional
geometry. Here a group is \emph{hyperbolic} if a Cayley graph for a
finite generating set has uniformly thin geodesic triangles; its boundary
consists of asymptotic classes of geodesic rays, with the usual visual
topology. An isometric action is \emph{geometric} if it is proper and
cocompact.

The following theorem resolves Cannon's conjecture positively.

\begin{theorem}[Cannon's conjecture]
\label{thm:main}
Let $G$ be a hyperbolic group with $\partial G$ homeomorphic to $\Sphere$.
There is a homomorphism $\rho:G\to\Isom(\HH^3)$ for which the action
on $\HH^3$ is proper and cocompact. Its kernel is finite.
\end{theorem}

Properness means that $\{g\in G:\rho(g)K\cap K\ne\varnothing\}$ is
finite for every compact $K\subset\HH^3$; cocompactness means that one
compact set has translates covering $\HH^3$. A finite kernel is allowed, and the image may contain orientation-reversing
isometries. In particular, the theorem does not require a choice of
orientation preserved by $G$.

For torsion-free $G$, the quotient is a closed hyperbolic
three-manifold, possibly nonorientable. We prove this consequence and
the resulting virtual-structure statements in \Cref{sec:consequences}.

\paragraph{History and the classical reduction.}
Cannon's 1994 combinatorial Riemann mapping theorem grew out of the
problem of finding an analytic structure on a sphere supplied only
with combinatorial boundary data. It recovers quasiconformal
coordinates from discrete annular moduli \cite[pp.~155--160]{Cannon}.
In their 1998 paper, Cannon--Swenson stated the group-boundary conjecture
and related a geometric action to conformality of the canonical boundary coverings
\cite[Theorem 5.1 and Conjecture 5.1]{CannonSwenson}.
Cannon--Floyd--Parry showed in 1999 that it suffices to find, around
each point in every prescribed neighborhood, an annulus whose discrete
moduli for separating curves stay bounded away from zero
\cite[Theorem 8.2]{CFP}.
Their finite subdivision rules provide a related framework for recursive
planar tilings \cite{CFPFiniteSubdivision}.
These works turn the topological problem into quantitative control of
crossings at small scales.

Bonk--Kleiner's uniformization theory gives a metric formulation of
this program \cite{BK}. We use the precise criterion in
Bourdon--Kleiner \cite[Corollary 3.5]{BourdonKleiner}: on an
approximately self-similar metric sphere, a uniform bound for the
combinatorial $2$-modulus of all continuous curves of a fixed sufficiently
small minimum diameter yields a quasi-M\"obius parametrization by the
round sphere. The visual boundary has the required self-similarity.
Thus the remaining analytic task is an upper modulus bound for its
fixed-scale coverings. The criterion and its hypotheses are verified
in \Cref{thm:uniformization}.

A related route uses \emph{Ahlfors regular conformal dimension}: the
infimum of the Hausdorff dimensions of Ahlfors regular metrics
quasisymmetric to the visual metric. This is a version of the
boundary invariant introduced by Pansu \cite{Pansu1989}.
Bonk--Kleiner proved Cannon's conjecture when that infimum is attained
\cite[Theorem 1.1]{BKConformalDimension}; their argument connects the
minimizing metric to quantitative curve families through the Loewner
property \cite[Theorem 1.3]{BKConformalDimension}. The existence of a
minimizing metric is a hypothesis of that criterion.

The algebraic approaches connect surface subgroups, cubulations, and
Kleinian realizations. Kahn--Markovi\'c constructed almost geodesic
surface subgroups in closed orientable hyperbolic three-manifold groups
\cite[Theorems 1.1--1.2]{KahnMarkovic}. Bergeron--Wise used this
abundance to cubulate those groups, and more generally proved that
quasiconvex codimension-one subgroups separating every distinct boundary
pair yield a proper cocompact cubical action
\cite[Theorem 1.4 and Section 5]{BergeronWise}.
Agol proved virtual specialness for hyperbolic groups with such an
action \cite[Theorem 1.1]{Agol}; these results give the consequences
recorded in \Cref{sec:consequences} once a hyperbolic three-manifold
realization is known.

In the recognition direction, Markovi\'c showed that enough
quasiconvex surface subgroups separating every boundary pair give a
Kleinian realization for groups whose sphere-boundary action is faithful
and orientation-preserving \cite[Theorem 1.1]{Markovic}.
Ha\"issinsky established a virtual Kleinian criterion for cubulated
hyperbolic groups with planar boundary \cite[Theorem 1.10]{Haissinsky}.
More recently, Groves, Ha\"issinsky, Manning, Osajda, Sisto, and Walsh
constructed group-theoretic drillings and showed that the Toral Relative
Cannon conjecture implies a virtual Kleinian realization for every
residually finite hyperbolic group with sphere boundary
\cite[Theorems A--B and Corollary C]{Drilling}.
The present proof follows the analytic route, establishing the required
modulus bound directly from hyperbolicity and the topology of the boundary.

After boundary uniformization, the invariant-conformal-structure method
of Sullivan and the detailed circumcenter construction of Tukia
straighten the action \cite{Sullivan,Tukia}.
We state the analytic inputs at their point of use and include the
short argument with unoriented conformal structures, so that
orientation-reversing elements are retained. Properness and
cocompactness then follow from the boundary action on ordered triples
(\Cref{prop:action}).

\paragraph{The proof in outline.}
Fix a visual metric on $Z=\partial G$. At scale $a^{-n}$, where $a>1$,
cover $Z$ by balls centered on a separated net. Assign nonnegative
weights to the balls so that every path of a fixed minimum diameter
meets balls of total weight at least one. Let $M(n)$ be the least sum
of squared weights. The classical criterion reduces the theorem to
\[
 \sup_n M(n)<\infty.
\]
Assume this fails. Record indices permit comparison of $M(N-l)$
with $M(N)$, and lead to two growth cases.

First, crossing duality in a coordinate rectangle produces weights
with squared sum $O(M(N)^{-1})$. A planar regularization turns them
into continuous functions $H_N$ on a slightly enlarged rectangle,
with separating level continua of uniformly positive diameter.
Annular estimates give a nonconstant locally uniform limit $u$ and
a finite measure $\mu$ such that
\[
 \bigl(\osc_{\overline B(x,s)}u\bigr)^2
 \le C\mu\bigl(\overline B(x,C_\mu s)\bigr).
\]
The approximants retain vanishing local squared oscillation sums
after every \emph{fixed} group transformation. A finite-cell
inequality compares two such functions by sweeping through the
separating levels. This is the analytic construction in
\Cref{sec:scalar}.

The growth of the record moduli determines how we use this scalar
limit. Exponential growth makes $u$ locally H\"older, which allows
normalized magnifications of its variation to converge to a nonconstant
function on a punctured sphere. The level comparison forces each transform with a different
puncture to be constant on the variation set of the first. A second
magnification contradicts this constraint (\Cref{sec:exponential}).

For subexponential growth, we use the measure bound to distribute
variation among many balls. A tree of separated balls carries nested
value intervals and a probability flow. The transition probabilities
obey linear and quadratic bounds in the ratio of child to parent
interval lengths. These bounds force a definite total probability,
at arbitrarily deep levels, onto nodes with two separated regions of variation
(\Cref{sec:flow}). At a common depth, many expansions have the same
two target regions but distinct punctures. Comparing them in pairs
charges more mass to small balls than their overlap allows
(\Cref{sec:pairs}). This excludes the remaining case.

Discrete crossing duality is
part of the extremal-length theory developed in
\cite{Schramm,CFPSquaring}. The finite incidence-vector identity is
also the probabilistic modulus formula of
Albin--Poggi-Corradini \cite[Theorem 5.1]{AlbinPoggiCorradini}; we
reprove it in the notation of our cell covers. Pansu's coarse conformal geometry uses
image-diameter energies and condenser capacities \cite[\S6]{Pansu1989}.
Scalar extremals and separating levels also enter Rajala's
uniformization of metric surfaces \cite[Sections 4--6]{RajalaUniformization}.
Here two path envelopes produce the required continua while preserving
finite-cell oscillation bounds. The subsequent level integrations and
pair counts are proved directly for the visual boundary. The two growth
branches use this same scalar construction, and their combination
finishes the theorem in \Cref{sec:completion}.

\paragraph{Conventions.}
A path is a continuous map of a compact interval, identified with its
image when convenient. A continuum is a
nonempty compact connected set; it need not be path connected.
$B(x,r)$ is open and $\overline B(x,r)=\{y:d(x,y)\le r\}$ is closed.
For real $f$, write $\osc_A f=\sup_A f-\inf_A f$.
Constants in $\lesssim$ and $\asymp$ are independent of the scale
indices being varied. Dependence on a compact region, a fixed group
element, or a subdivision parameter is specified where needed.

\section{Boundary geometry and the modulus criterion}
\label{sec:geometry}

We collect the geometric estimates used throughout the proof, then state
the classical uniformization input. Let $\Gamma$ be the Cayley graph of
$G$ for a finite symmetric generating set, with unit edge lengths. Write
$e$ for its identity vertex and $|v|=d_\Gamma(e,v)$.
The Gromov product is
\[
 (y|z)_o=\tfrac12\bigl(d_\Gamma(o,y)+d_\Gamma(o,z)-d_\Gamma(y,z)\bigr).
\]
Boundary products are defined by limits along rays, up to a uniformly
bounded additive error. A visual metric $d$ on $Z=\partial G$ satisfies
\begin{equation}
 C_v^{-1}a^{-(y|z)_e}\le d(y,z)\le C_va^{-(y|z)_e},
 \qquad y\ne z,
 \label{eq:visual}
\end{equation}
for some $a>1$. Set $r_k=a^{-k}$ for integers $k\ge0$.
We use the standard product inequality
$(y|z)_e\ge\min\{(y|w)_e,(w|z)_e\}-O(1)$ and its geometric form:
geodesics with product at least $j-C$ stay a bounded synchronous
distance apart through depth $j$, provided both have that length.
The bound depends on $C$ and the
hyperbolicity constant. These facts, visual metrics, and thin ideal
triangles are recalled in \cite[III.H.3]{BH}.

\subsection{Expansions}

Left multiplication by the inverse of a vertex far along a ray magnifies
a small boundary ball. The following formula also controls the complement
of that ball; this second feature will be essential.

\begin{lemma}[Expansion formula]
\label{lem:expansion}
Let $v$ be the depth-$k$ vertex on a based ray to $x\in Z$, put
$h=v^{-1}$ and $R=r_k$, and set $D_y=\max\{R,d(x,y)\}$.
There are constants $c_E,C_E>0$, independent of these choices, such that
\begin{equation}
 c_E\frac{R\,d(y,z)}{D_yD_z}
 \le d(hy,hz)\le C_E\frac{R\,d(y,z)}{D_yD_z}.
 \label{eq:expansion}
\end{equation}
In particular, on $B(x,CR)$ the map $h$ is bi-Lipschitz with scale
factor $R^{-1}$ and distortion depending only on $C$.
If $d(x,z_0)\ge\varepsilon_*>0$, $R<\varepsilon_*/2$, and $p=hz_0$, then
\begin{equation}
 d(hy,p)\lesssim R/D_y\quad(y\in Z),\qquad
 d(hy,p)\asymp R/D_y\quad(d(x,y)\le\varepsilon_*/2),
 \label{eq:puncture}
\end{equation}
and
\begin{equation}
 (v^{-1}|p)_e\ge k-C(\varepsilon_*).
 \label{eq:inverse-prefix}
\end{equation}
The constants in the last two estimates may depend on $\varepsilon_*$.
\end{lemma}
\begin{proof}
Base change and thin triangles give
\[
 (y|z)_v=(y|z)_e+k-(v|y)_e-(v|z)_e+O(1),
 \qquad (v|y)_e=\min\{k,(x|y)_e\}+O(1).
\]
Thus $a^{-(v|y)_e}\asymp D_y$. Exponentiating the first identity
and using $(hy|hz)_e=(y|z)_v+O(1)$ proves
\eqref{eq:expansion}. With $z=z_0$, the factor $D_{z_0}$ is bounded
above and below, and $d(y,z_0)$ is bounded above. When
$d(x,y)\le\varepsilon_*/2$, the latter is also bounded below.
This proves \eqref{eq:puncture}. Finally,
$(v^{-1}|v^{-1}z_0)_e=k-(v|z_0)_e+O(1)$, and
$(v|z_0)_e$ is bounded in terms of $\varepsilon_*$.
\end{proof}

\begin{lemma}[Prefixes]
\label{lem:prefix}
Every vertex $v$ lies a uniformly bounded distance from a depth-$|v|$
vertex on a based ray. With that ray's endpoint as $x$,
\eqref{eq:expansion} remains valid for $h=v^{-1}$ and $R=r_{|v|}$,
with uniform constants. Moreover:
\begin{enumerate}[label=\textup{(\roman*)}]
\item If $d(\xi,\eta)\le Cr_j$, depth-$j$ vertices on any based
rays to $\xi,\eta$ are a bounded distance apart, depending only on $C$.
\item If $|v|=k$ and $(v|\xi)_e\ge k-C$, every geodesic from $e$
to $v$ stays a bounded synchronous distance from a ray to $\xi$
through depth $k$.
\end{enumerate}
\end{lemma}
\begin{proof}
Translate a bi-infinite geodesic so that it passes through $v$.
The ideal triangle with third vertex $e$ puts $v$ uniformly near one
of its based sides; the nearby parameter differs from $|v|$ by a
bounded amount. A bounded change of basepoint changes products by a
bounded amount, preserving \eqref{eq:expansion}. The two remaining
assertions follow from \eqref{eq:visual} and the product estimate.
\end{proof}

Each fixed $g\in G$ acts bi-Lipschitzly on $Z$, with a constant
depending on $g$. The whole action is uniformly quasi-M\"obius:
base-change factors cancel in distance cross-ratios, giving a single
$C\ge1$ such that
\[
 C^{-1}\frac{d(x,z)d(y,w)}{d(x,w)d(y,z)}
 \le\frac{d(gx,gz)d(gy,gw)}{d(gx,gw)d(gy,gz)}
 \le C\frac{d(x,z)d(y,w)}{d(x,w)d(y,z)}
\]
for every distinct quadruple and every $g$.

\subsection{Packing and paths}

\begin{lemma}[Packing]
\label{lem:packing}
There are $Q\ge2$ and $C_P<\infty$ such that an $s$-separated subset
of a ball of radius $t$ has at most $C_P(t/s)^Q$ points whenever
$0<s\le t\le\diam Z$. In particular, $Z$ is doubling.
\end{lemma}
\begin{proof}
Choose $k$ with $r_k\asymp t$. Rays ending in the given ball have
their depth-$k$ vertices in a uniformly bounded graph ball, by
\Cref{lem:prefix}. At depth $k+j$ there are at most $CA^j$ possible
continuations, where $A>1$ depends on the valence. Endpoints whose
rays share such a vertex are at visual distance at most $C'r_{k+j}$.
Choose the least $j\ge0$ with $C'r_{k+j}<s$. The desired bound
follows with $Q\ge\max\{2,\log_a A\}$. Coarse scales are absorbed
by increasing the constant.
\end{proof}

\Needspace{5\baselineskip}
\begin{lemma}[Controlled paths]
\label{lem:paths}
There is $L\ge1$ such that any two points of $\overline B(x,s)$
join by a path in $B(x,Ls)$. Any two points of a continuum can also
be joined by a path in any prescribed open neighborhood of it.
\end{lemma}
\begin{proof}
For small $s$, expand at $x$ with $R\asymp s$ and choose $z_0$ a
uniformly positive distance from $x$; two fixed distinct points
suffice for this choice. By \eqref{eq:puncture}, the two image
points stay a uniform positive distance from $p=hz_0$.
Fix a homeomorphism $\phi:Z\to\Sphere$. Uniform continuity of
$\phi^{-1}$ places their spherical images a uniform distance from
$\phi(p)$. Join them outside a smaller round ball about $\phi(p)$.
Uniform continuity of $\phi$ makes the pullback path stay a uniform
$d$-distance from $p$. The upper bound in \eqref{eq:puncture}
puts its preimage under $h$ in $B(x,Ls)$. Increase $L$ for the
remaining scales, using path connectedness of $Z$.

For a continuum $C$ in an open set $V$, choose $\varepsilon>0$
with its $L\varepsilon$-neighborhood in $V$. Connectedness supplies
a finite $\varepsilon$-chain in $C$ between any two specified
points: the set reachable by such chains is relatively clopen.
Join consecutive points by the first assertion.
\end{proof}

\subsection{The classical reduction}

Choose a maximal $r_n$-separated set $T_n\subset Z$ and its cells
$B(q,2r_n)$, $q\in T_n$. These cover $Z$; fixed enlargements have
bounded overlap. For a path family $\mathcal F$, define
\[
 \Mod_2(\mathcal F,T_n)=\min_{\omega\ge0}
 \left\{\sum_{q\in T_n}\omega(q)^2:
 \sum_{B(q,2r_n)\cap\pi\ne\varnothing}\omega(q)\ge1
 \text{ for every }\pi\in\mathcal F\right\}.
\]
Only finitely many incidence vectors occur, so the minimum exists.
Fix a sufficiently small $d_0>0$ and put
\begin{equation}
 \mathcal F_0=\{\pi:\diam\pi\ge d_0\},\qquad
 M(n)=\Mod_2(\mathcal F_0,T_n).
 \label{eq:macro-modulus}
\end{equation}
We require $d_0<\diam Z/10$ and choose it so that every path of
diameter at least $R/2$ in $\overline B(x,3R)$ has image diameter
at least $d_0$ under a scale-$R$ expansion. This is possible by
\eqref{eq:expansion}. We also impose the smallness requirement in
\Cref{thm:uniformization} before fixing $d_0$ and $M(n)$.
The family is nonempty, and $0<M(n)<\infty$: all-one weights are
admissible, while any one path gives the lower bound $1/\#T_n$.

\begin{theorem}[Bonk--Kleiner; Bourdon--Kleiner]
\label{thm:uniformization}
If $\sup_n M(n)<\infty$, then $(Z,d)$ is quasi-M\"obius
homeomorphic to the round sphere.
\end{theorem}
\begin{proof}[Verification of the cited criterion]
Corollary 3.5 of \cite{BourdonKleiner}, derived from
\cite[Theorem 10.4]{BK}, applies to an arcwise connected,
approximately self-similar metric sphere. Approximate self-similarity
means that each open ball, with distances divided by its radius,
is uniformly bi-Lipschitz to an open subset of the space.
\Cref{lem:expansion} supplies this property; intermediate and large
radii change only the constant.

The criterion assumes bounded $2$-moduli of paths of diameter at
least a sufficiently small $d_0$ in successive graph approximations.
For a dyadic scale $2^{-m}$ choose $n$ with
$2^{-m}/a<r_n\le2^{-m}$. Our cells form a fixed $\kappa$-approximation
for any $\kappa>2a$: the inner balls of radius $2^{-m}/\kappa$
are disjoint, each is contained in its cell, and each cell is
contained in the radius-$\kappa2^{-m}$ ball about its center.
The modulus is exactly $M(n)$. Finitely many coarse scales have
finite modulus. No Coxeter-group assumption occurs in this criterion.
\end{proof}

\begin{proposition}[From the boundary to hyperbolic space]
\label{prop:action}
If $(Z,d)$ is quasi-M\"obius homeomorphic to the round sphere,
then $G$ admits a proper cocompact isometric action on $\HH^3$
with finite kernel.
\end{proposition}
\begin{proof}
Let $\psi:Z\to\Sphere$ be the uniformizing homeomorphism and write
$\mathcal G=\{\psi g\psi^{-1}:g\in G\}$. The boundary cross-ratio estimate
and the two fixed maps $\psi,\psi^{-1}$ give one quasi-M\"obius distortion
function for every element of $\mathcal G$. Fix $f\in\mathcal G$ and $x\in\Sphere$,
keep a fourth point away from $x$, and let two points at equal spherical
distance from $x$ approach $x$. In the cross-ratio inequality the factors
involving the fourth point tend to one, both before and after applying
$f$. Thus the infinitesimal linear distortion has a bound independent
of $f$ and $x$; the radius at which the estimate is accurate need not
be uniform in $f$. The metric characterization of planar
quasiconformality, applied in conformal charts, makes each $f$
quasiconformal. At its almost-everywhere nonsingular differentiability
points, the same infinitesimal bound controls the ratio of singular values
of $Df$, and hence its analytic dilatation. This gives a common
quasiconformal bound. For a reversing map apply the characterization
after a conformal reflection; see \cite[Theorems 2.5.3--2.5.4]{AIM}
and the metric characterization in \cite[Theorem 1.4 and Remark 1.8(b)]{HK}.

We recall explicitly the invariant-structure argument, including
orientation reversal. This is the invariant-conformal-structure method of
Sullivan, with the circumcenter construction detailed by Tukia
\cite[pp.~73--76]{Tukia}. Analytic quasiconformal maps have
almost-everywhere defined nonsingular weak differentials, and they and
their inverses preserve null sets \cite[Theorem 3.1.2]{AIM}.
Their differentials satisfy the almost-everywhere composition rule
\cite[Lemma 10 and Theorem 5]{AhlforsBers}.
Since $\mathcal G$ is countable, remove the $\mathcal G$-orbit
of the union of all exceptional sets for these assertions and for all
pairs of maps. Its complement $X$ is invariant and conull. At $x\in X$,
let $\mathcal C_x$ be the space of positive definite quadratic forms on
$T_x\Sphere$ modulo positive scalars. This is a hyperbolic plane, and
pullback by any invertible linear map, of either orientation, is an
isometry between the corresponding fibers
\cite[Section 2.3, equations (2.17)--(2.18) and Theorem 2.3.2]{AIM}. The orbit
$\{f^*q_{\rm round}(x):f\in\mathcal G\}\subset\mathcal C_x$ is uniformly
bounded, by the common quasiconformal bound, and contains the round
form. Its unique hyperbolic circumcenter defines a uniformly bounded
measurable conformal structure $q(x)$. Indeed, circumcenters of the
increasing finite subsets from an enumeration are measurable and
converge to that circumcenter: their radii increase, their centers stay
in a fixed compact ball, and every subsequential limit is the unique
center of the whole orbit. The composition rule permutes this orbit
under pullback, so $f^*q=q$ on $X$ for every $f\in\mathcal G$. Extend $q$
arbitrarily across the null complement of $X$.

The measurable Riemann mapping theorem
\cite[Theorem 6]{AhlforsBers} gives a quasiconformal
homeomorphism $F:\Sphere\to\Sphere$ with $F^*q_{\rm round}=q$ almost
everywhere. Invariance and the analytic composition rule show that
every $FfF^{-1}$ is conformal or anticonformal, according to the
orientation of $f$. Such sphere homeomorphisms are M\"obius or
anti-M\"obius. Their unique extensions to hyperbolic isometries therefore
define a homomorphism $\rho:G\to\Isom(\HH^3)$.

For completeness, properness and cocompactness follow directly
from the action on ordered distinct triples $\mathcal T(Z)$.
On any compact set of triples all pairwise distances are bounded
below. For $|v|\to\infty$, at least two points of such a triple
stay away from the expansion center for $v^{-1}$; their image
distance is $O(r_{|v|})$. Thus $v^{-1}K\cap K=\varnothing$
eventually for every compact $K\subset\mathcal T(Z)$.
Conversely, expand at one point of a closest pair with scale
comparable to their distance. Formula \eqref{eq:expansion} makes
all three image distances uniformly positive: if the third point
is nearby, all denominator factors have that scale; otherwise
its distances to the pair cancel its denominator factor.
Triples already uniformly separated need no expansion.
Hence the action on triples is proper and cocompact.

These properties survive boundary conjugacy. Let $I=\Isom(\HH^3)$.
The stabilizer $H$ of a standard ordered boundary triple has two
elements: the identity and reflection in its hyperbolic plane.
Choose $o\in\HH^3$ fixed by $H$, and let $K_o$ be its compact
stabilizer in $I$. The equivariant map
\[
 \mathcal T(\Sphere)\cong I/H\longrightarrow I/K_o\cong\HH^3
\]
is continuous, onto, and proper: isometries sending $o$ into a
compact set form a compact set. Properness and cocompactness
therefore descend. The kernel fixes $o$, so is finite.
\end{proof}

\section{A scalar limit from unbounded modulus}
\label{sec:scalar}

Assume that the macroscopic moduli $M(n)$ are unbounded. Crossing
duality first gives transverse weights of energy $O(M(n)^{-1})$ along
a subsequence. We turn those weights into continuous functions whose
level continua meet every path straddling the corresponding value.
An annular estimate then yields a nonconstant limit $u$ and a finite
measure controlling its squared oscillation. The approximating functions
retain vanishing local oscillation energy after each fixed group
transformation; later comparisons will use this property as well as the
limit itself. Constants are independent of the approximation index;
thresholds may depend on the fixed compact subset of the domain under
consideration.

\subsection{Record scales and crossing duality}

If $\limsup M(n)^{1/n}>1$, choose
$1<\lambda<\limsup M(n)^{1/n}$; otherwise put $\lambda=1$.
In both cases $M(n)/\lambda^n$ is unbounded. Retain increasing record
indices $N$ whose record values tend to infinity. Thus $M(N)\to\infty$ and
\begin{equation}
 M(N-l)\leq\lambda^{-l}M(N)\qquad(0\leq l\leq N).
 \label{eq:records}
\end{equation}
This inequality survives every further subsequence.

\begin{lemma}[Localization to a rectangle]
\label{lem:rectangle-localization}
There are coordinate rectangles
$\cQ=[0,1]^2\subset\Qhat=[0,1]\times[-1,2]\subset Z$ and a
subsequence of the record indices such that the left-to-right crossing
modulus of $\cQ$ is at least $cM(N)$, with $c>0$ independent of $N$.
\end{lemma}
\begin{proof}
Cover $Z$ by finitely many coordinate boxes $[-1,1]^2$, each lying in
a chart containing $[-8,8]^2$ of diameter less than $d_0$. A macroscopic
path meeting an inner box must leave its larger box. Follow it to its
first meeting with $\partial[-2,2]^2$. If it exits through $x=2$,
its segment after the last preceding meeting with $x=1$ crosses
$[1,2]\times[-2,2]$ from left to right. The other three exit sides
give analogous strips. Each strip extends by its height in both
transverse directions inside $[-8,8]^2$.

Consequently there are finitely many crossing families
$\mathcal F_1,\ldots,\mathcal F_J$ such that every macroscopic path
contains a member of one of them. Summing admissible weights gives
\[
 M(N)\leq J\sum_{j=1}^J\Mod_2(\mathcal F_j,T_N)
       \leq J^2\max_j\Mod_2(\mathcal F_j,T_N).
\]
One strip realizes this bound along a subsequence. Affine coordinates
give the asserted rectangles.
\end{proof}

The edge $y=0$ is the top of $\cQ$. The rectangles
$[0,1]\times[-1,0]$ and $[0,1]\times[1,2]$ are its top and bottom
bands, respectively. Coordinates serve only interpolation and topology:
all distances refer to $d$.

\begin{lemma}[Crossings and blockers]
\label{lem:planar-blocker}
In a closed coordinate rectangle:
\begin{enumerate}[label=\textup{(\roman*)}]
\item every continuum meeting both lateral edges meets every
top-to-bottom path;
\item every compact set meeting all top-to-bottom paths contains
a continuum meeting both lateral edges.
\end{enumerate}
\end{lemma}
\begin{proof}
For (i), suppose a lateral continuum and a top-to-bottom path are
disjoint. Enlarge the rectangle, extending the continuum horizontally
and the path vertically to the corresponding new edges. The extensions
remain disjoint: a new horizontal and vertical segment can meet only at
an original corner, already a common point of the original sets.
Approximate the extended path, within its positive clearance from the
continuum, by a polygonal path with interior in the enlarged rectangle.
After subdivision at intersections it contains a simple polygonal arc
between the top and bottom edges. The Jordan curve theorem, applied
after adjoining a boundary route, shows that this arc separates the
lateral edges. The connected continuum cannot avoid it.

For (ii), subdivide into finer square grids and mark the closed squares
meeting the blocker $K$. No edge-adjacent chain of unmarked squares
joins top to bottom. Label $A$ the unmarked squares reachable from the
top through unmarked edge-adjacent squares, and append a top row of
$A$ squares. Label all remaining squares and an appended bottom row $B$.
Draw the interfaces between unlike labels. Interior vertices have even
degree; each lateral boundary has an odd number of ends, since its
labels begin with $A$ and end with $B$. There are no top or bottom ends.
Some interface component therefore joins the lateral boundaries:
otherwise every component contributes an even number of ends to the
only lateral boundary it meets.

Along a path in this component, the adjacent $B$ squares are original
marked squares. An adjacent unmarked square would be reachable, while
an interface with the appended bottom row would give an unmarked
top-to-bottom chain. Successive $B$ squares meet at least at a corner.
Their union is a connected compact lateral crossing. A Hausdorff limit
as the mesh tends to zero is connected, meets both lateral edges, and
lies in $K$.
\end{proof}

The next finite-dimensional argument belongs to discrete extremal
length. Cannon--Floyd--Parry describe optimal weights through crossing
vectors and their pairing with transverse cuts
\cite[Sections 2.3 and 2.4.3]{CFPSquaring}; the probability formulation
of Albin--Poggi-Corradini gives the same least-norm incidence-vector
identity \cite[Theorem 5.1]{AlbinPoggiCorradini}. We give the argument
for the cell families used here.

\begin{lemma}[Crossing dual weights]
\label{lem:dual-weights}
For each retained $N$ there are nonnegative weights
$\rho_q=\rho_q^{(N)}$, $q\in T_N$, admissible for top-to-bottom
crossings of $\cQ$, such that
\begin{equation}
 \sum_{q\in T_N}\rho_q^2\leq \frac{C}{M(N)}.
 \label{eq:dual-energy}
\end{equation}
\end{lemma}
\begin{proof}
Let $P$ minimize Euclidean norm in the convex hull of the finitely many
cell-incidence vectors $v$ of left-to-right crossings. Its coordinates
are nonnegative and sum to at least one, so $P\ne0$. Minimality gives
$\langle P,v\rangle\geq\|P\|^2$, making $P/\|P\|^2$ admissible.
Conversely any admissible $\omega$ satisfies
$\langle\omega,P\rangle\geq1$. Thus the left-to-right modulus is
$\|P\|^{-2}$. By \Cref{lem:planar-blocker}(i), each top-to-bottom
incidence vector has inner product at least one with every $v$, hence
with $P$. Set $\rho=P$ and apply \Cref{lem:rectangle-localization}.
\end{proof}

\subsection{Continuous functions with separating levels}

For $\pi\subset\Qhat$ write
$\ell_\rho(\pi)=\sum_{q:B(q,2r_N)\cap\pi\ne\varnothing}\rho_q$.
For $x\in\cQ$, let $b_N(x)$ be the infimum of this cost over paths
from the top edge to $x$ in $\cQ$. Define the raw function
$\min\{1,b_N\}$, assigning exactly zero on the top edge, and extend
it by zero and one to the top and bottom bands. Admissibility gives
value one on the bottom edge.

Along every path in $\Qhat$, the difference of raw endpoint values is
at most its cost. Inside $\cQ$ this follows by concatenating with
nearly minimizing paths and clipping at one; a path beginning on the
top edge is itself a competitor. A path visiting both bands contains
a crossing of cost at least one. If it visits only the top band, each
endpoint value is bounded by the cost of its segment to that band.
If it visits only the bottom band, each endpoint value is at least one
minus that segment's cost, by concatenation with a top-to-endpoint
competitor. In either case the endpoint difference is bounded by the
cost of the entire path.

Triangulate $\Qhat$ by a finite coordinate grid containing the band
boundaries, with every triangle of $d$-diameter less than $r_N$.
Interpolate the raw vertex values affinely to obtain a continuous
$U_N^0:\Qhat\to[0,1]$. It is constant on the bands. Continuity alone
does not ensure that one selected level continuum meets every
straddling path. The following two path envelopes give the strict
sublevel and superlevel sets the connections to the bands needed for
that conclusion. Define
\[
 F_N(x)=\inf_{\pi:x\to\mathrm{top}}\max_\pi U_N^0,
 \qquad
 H_N(x)=\sup_{\pi:x\to\mathrm{bottom}}\inf_\pi F_N,
\]
where paths remain in $\Qhat$ and end in the indicated band.

\begin{lemma}[Regularization and level continua]
\label{lem:level-continua}
The functions $F_N,H_N$ are continuous, with the prescribed band values,
and for every path $K\subset\Qhat$,
\begin{equation}
 \osc_K H_N\leq\osc_K F_N\leq\osc_K U_N^0.
 \label{eq:flood-osc}
\end{equation}
For every $0<t<1$, the set $\{H_N=t\}\cap\cQ$ contains a lateral
crossing continuum $C_{N,t}$ of diameter at least a fixed $d_1>0$.
Every path in $\Qhat$ whose endpoint values strictly straddle $t$
meets $C_{N,t}$.
\end{lemma}
\begin{proof}
We have $F_N\geq U_N^0$ and $H_N\leq F_N$; constant paths give
the band values. Concatenation with nearly optimizing paths gives,
for $x,y\in K$,
\[
 F_N(x)\leq\max\{\max_K U_N^0,F_N(y)\},\qquad
 H_N(y)\geq\min\{\inf_K F_N,H_N(x)\}.
\]
If $F_N(x)>F_N(y)$, the first inequality places their values between
$\min_K U_N^0$ and $\max_K U_N^0$. If $H_N(x)>H_N(y)$, the
second places them between $\inf_K F_N$ and $\sup_K F_N$.
This proves \eqref{eq:flood-osc}. Short coordinate segments and
continuity of $U_N^0$ then prove continuity of both functions,
including at the boundary.

Every point of $\{H_N>t\}$ joins the bottom band within that set:
choose a path with $\inf F_N>t$ and use its tails. Likewise every
point of $\{F_N<t\}$ joins the top within $\{F_N<t\}$.
To connect a point $x$ with $H_N(x)<t\leq F_N(x)$ to the top,
choose $H_N(x)<s<t$ and follow any top-going path until its first
point with $F_N=s$. On this prefix $F_N\geq s$, so concatenation
with nearly maximizing paths from any prefix point $z$ gives
\[
 H_N(x)\geq\min\{s,H_N(z)\}.
\]
Hence $H_N(z)\leq H_N(x)<t$ throughout the prefix. Its endpoint
joins the top through $F_N<t$, hence through $H_N<t$. Thus both
strict sides of the level join their respective bands, and then the
outer top and bottom edges, without meeting the level.

The compact set $\{H_N=t\}$ blocks all top-to-bottom paths by
continuity and lies in $\cQ$. By \Cref{lem:planar-blocker}(ii) it
contains a lateral continuum $C_{N,t}$; take $d_1$ to be the distance
between the lateral edges of $\cQ$. Complete any straddling path to
a top-to-bottom path using the strict-side connections. The completed
path meets $C_{N,t}$ by \Cref{lem:planar-blocker}(i), and the added
pieces avoid the level, proving the last assertion.
\end{proof}

Figure~\ref{fig:level-continuum} isolates the last step: the strict-side
extensions cannot account for the intersection with the selected
continuum, so the original straddling path must meet it.

\begin{figure}[htbp]
\centering
\begin{tikzpicture}[x=1cm,y=1cm,font=\small,
  line cap=round,line join=round,>=Stealth]
  \definecolor{levelblue}{RGB}{54,101,125}
  \definecolor{levelochre}{RGB}{164,111,48}
  \fill[black!3] (0,1.25) rectangle (7,1.72);
  \fill[black!3] (0,-1.72) rectangle (7,-1.25);
  \draw[black!45] (0,-1.72) rectangle (7,1.72);
  \draw[black!45] (0,1.25)--(7,1.25) (0,-1.25)--(7,-1.25);
  \node at (1.9,1.49) {top band: $H_N=0$};
  \node at (1.9,-1.49) {bottom band: $H_N=1$};
  \node[anchor=north west,text=black!65] at (.10,1.15) {$\mathcal Q$};
  \draw[levelblue,line width=2pt]
    (0,.13)--(.37,.03)--(.69,.19)--(1.04,.05)--(1.34,.20)
    --(1.64,-.04)--(1.98,.08)--(2.28,-.09)--(2.60,.04)
    --(2.91,-.13)--(3.19,.06)--(3.53,-.09)--(3.84,.09)
    --(4.16,-.02)--(4.48,.15)--(4.81,.01)--(5.13,.18)
    --(5.47,.04)--(5.80,.20)--(6.10,.05)--(6.41,.20)
    --(6.70,.08)--(7,.20);
  \draw[levelblue,line width=1.1pt]
    (1.34,.20)--(1.42,.37)--(1.59,.32)
    (2.28,-.09)--(2.18,-.28)--(2.38,-.24)
    (4.81,.01)--(4.89,-.21)--(5.03,-.15)
    (6.10,.05)--(6.04,-.14)--(6.22,-.19);
  \node[levelblue,anchor=south] at (1.95,.40) {$C_{N,t}\subset\{H_N=t\}$};
  \draw[levelochre,dashed,thick]
    (4.9,1.72) .. controls (4.4,1.32) and (3.88,1.29) .. (3.72,.79);
  \draw[levelochre,thick]
    (3.72,.79) .. controls (3.50,.34) and (4.24,-.24) .. (4.10,-.80);
  \draw[levelochre,dashed,thick]
    (4.10,-.80) .. controls (4.00,-1.24) and (4.94,-1.29) .. (5.1,-1.72);
  \fill[levelochre] (3.72,.79) circle (1.5pt) (4.10,-.80) circle (1.5pt);
  \node[anchor=west] at (3.96,.74) {$H_N<t$};
  \node[anchor=east] at (3.89,-.86) {$H_N>t$};
  \node[levelochre,anchor=west] at (4.27,-.46) {$\eta$};
  \node[anchor=west,text=black!85,align=left] at (7.32,.82)
    {strict-side\\extension};
  \draw[black!65,thin] (7.2,.86)--(5.10,1.18);
  \node[anchor=west,text=black!85,align=left] at (7.32,-.85)
    {strict-side\\extension};
  \draw[black!65,thin] (7.2,-.82)--(5.05,-1.36);
\end{tikzpicture}
\caption{A separating level continuum and a straddling path. The dashed
extensions stay in the strict sublevel and superlevel sets and complete
$\eta$ to a top-to-bottom crossing. That crossing meets $C_{N,t}$, so
$\eta$ itself meets $C_{N,t}$. The continuum is drawn schematically;
neither smoothness nor path connectedness is assumed.}
\label{fig:level-continuum}
\end{figure}


\subsection{Oscillation and annular estimates}

The level continua now let us estimate oscillation by crossing an annulus.
Energies formed from image diameters and their relation to condenser
capacity appear in Pansu's coarse conformal geometry \cite[\S6]{Pansu1989};
the estimates below use the finite-cell energy of our regularized scalar
functions. Put $U=\operatorname{int}_Z\Qhat$. For a continuous function $f$ on
an open set containing a compact set $K$, and $J\geq2$, define
\[
 \mathcal E_n^{(J)}(f;K)=
 \sum_{\substack{q\in T_n\\B(q,2r_n)\cap K\ne\varnothing}}
       \bigl(\osc_{\overline B(q,Jr_n)}f\bigr)^2.
\]
The closed balls lie in the domain for all sufficiently large $n$.

\begin{lemma}[Local oscillation]
\label{lem:local-oscillation}
For fixed compact $K\subset U$ and $J>0$, uniformly for $x\in K$
and sufficiently large $N$,
\begin{equation}
 \osc_{\overline B(x,Jr_N)}H_N
 \leq\sum_{\substack{q\in T_N\\d(q,x)\leq C(J)r_N}}\rho_q.
 \label{eq:local-osc}
\end{equation}
Here $C(J)$ depends only on $J$ and the ball path constant.
Consequently, for every fixed $g\in G$, compact $K'\subset gU$,
and $J\geq2$,
\[
 \mathcal E_N^{(J)}(H_N\circ g^{-1};K')
       \leq\frac{C(g,K',J)}{M(N)}\longrightarrow0.
\]
\end{lemma}
\begin{proof}
Join two points of $\overline B(x,Jr_N)$ inside $B(x,LJr_N)$
by \Cref{lem:paths}. Vertices of triangles meeting this path lie
in $B(x,(LJ+1)r_N)$; any two such vertices join inside
$B(x,L(LJ+1)r_N)\subset U$ for large $N$. The raw endpoint
inequality bounds every difference of their values by the right side
of \eqref{eq:local-osc}, with $C(J)=L(LJ+1)+3$. It therefore bounds
the interpolant's oscillation along the original path. Apply
\eqref{eq:flood-osc}.

Squaring and summing costs a bounded factor by \Cref{lem:packing}:
both the number of nearby weights and the number of centers charging
each weight are bounded. The same holds for $br_N$-separated centers,
with constants depending on $b>0$. A fixed $g$ is bi-Lipschitz,
so inverse images of target balls lie in source balls of radius
$C_gJr_N$, their centers are $b_gr_N$-separated, and all these source
balls lie in a fixed compact subset of $U$. The summation bound and
\eqref{eq:dual-energy} finish the proof.
\end{proof}

The mesh-scale estimate preserves the small dual energy. We next
compare oscillation across a ball of radius $R=r_l$ with the dual
energy in a surrounding annulus. Expansion takes that annulus to a
fixed-size region and takes the mesh $r_N$ to $r_{N-l}$; this is why
the next bound contains $M(N-l)$.

\begin{lemma}[Annular estimate]
\label{lem:annular}
For every compact $K\subset U$ there are $R_*>0$ and an integer
$m_0$ such that, if $x\in K$, $R=r_l\leq R_*$, and $N-l\geq m_0$,
then any two points joined by a path in $B(x,R/2)$ have
$H_N$-value difference $\Delta$ satisfying
\begin{equation}
 \Delta^2\leq C M(N-l)
       \sum_{\substack{q\in T_N\\R/4\leq d(q,x)\leq4R}}\rho_q^2.
 \label{eq:annular}
\end{equation}
The constant is uniform; $N$ is sufficiently large for the fixed
interior region.
\end{lemma}
\begin{proof}
Choose $R_*$ with $4R_*<d_1$ and all closed $5R_*$-balls about
$K$ contained in $U$. All following restrictions on $m_0$ concern
only $r_N/R=a^{-(N-l)}$.

Expand at $x,l$ by $h$, and choose macroscopic admissible weights
$\omega$ at scale $N-l$ with squared mass $M(N-l)$. Put
\[
 \mathcal A=\{q\in T_N:R/2\leq d(q,x)\leq3R\},\qquad
 \sigma_q=\sum_{\substack{\xi\in T_{N-l}\\
               d(\xi,hq)\leq C_0r_{N-l}}}\omega_\xi,
 \quad q\in\mathcal A,
\]
where $C_0>C_E+2$. Expansion on $B(x,3R)$ separates the $hq$ by
a fixed multiple of $r_{N-l}$. Packing bounds both overlap
multiplicities, giving
\[
 \sum_{q\in\mathcal A}\sigma_q^2\leq C M(N-l).
\]

Let $I$ be the open interval between the endpoint values. For $t\in I$,
the continuum $C_{N,t}$ meets the given path and exits $B(x,2R)$,
since its diameter exceeds $4R$. By \Cref{lem:paths}, join these
two continuum points in its $r_N/4$-neighborhood and retain a radial
subpath $\eta_t$ between radii $R$ and $2R$, contained in that closed
annulus. Its diameter is at least $R$, so $h\eta_t$ is macroscopic
by the choice of $d_0$ in \Cref{sec:geometry}.

If a target cell $B(\xi,2r_{N-l})$ meets $h\eta_t$ at $hy$,
choose $q\in T_N$ with $d(q,y)\leq r_N$. For large $m_0$, this
center belongs to $\mathcal A$ and
$d(hq,hy)\leq C_Er_N/R$, so $\omega_\xi$ occurs in $\sigma_q$.
Some level point lies within $r_N/4$ of $y$, hence in
$\overline B(q,2r_N)$. Thus $t$ lies in the closed range interval
$I_q$ of $H_N$ on that ball. Accounting for all target cells gives
\[
 1\leq\sum_{q\in\mathcal A}\sigma_q\mathbf1_{I_q}(t)
 \quad(t\in I),
 \qquad
 \Delta\leq C M(N-l)^{1/2}
       \left(\sum_{q\in\mathcal A}|I_q|^2\right)^{1/2}.
\]
The second inequality follows by integration and Cauchy--Schwarz.
Apply \eqref{eq:local-osc} with $J=2$ and its bounded-overlap
summation. Taking $C(2)r_N<R/4$ confines every charged center to
$R/4\leq d(q,x)\leq4R$, proving \eqref{eq:annular}.
\end{proof}

\subsection{The limiting function and measure}

The finite measures
$\mu_N=M(N)\sum_{q\in T_N}\rho_q^2\delta_q$
have uniformly bounded mass by \eqref{eq:dual-energy}.

\Needspace{5\baselineskip}
\begin{proposition}[Scalar limit]
\label{prop:scalar-limit}
Along a subsequence of the record indices,
$H_N\to u$ locally uniformly on $U$ and $\mu_N\rightharpoonup\mu$,
where $u:U\to[0,1]$ is continuous and nonconstant and $\mu$ is finite.
The function $u$ is zero and one on the interiors of the respective
bands, and is locally H\"older continuous if $\lambda>1$.
For every compact $K\subset U$ and sufficiently small
$s>0$, uniformly for $x\in K$,
\begin{equation}
 \bigl(\osc_{\overline B(x,s)}u\bigr)^2
       \leq C\mu\bigl(\overline B(x,C_\mu s)\bigr).
 \label{eq:measure}
\end{equation}
For each $0<t<1$ there is a continuum $C_t\subset\cQ$ of diameter
at least $d_1$, with $u=t$ on $C_t\cap U$, meeting every compact
path in $U$ whose endpoint values strictly straddle $t$.
The fixed-transform energy bounds of \Cref{lem:local-oscillation}
hold along the retained subsequence.
\end{proposition}
\begin{proof}
Fix compact $K\subset U$, and let $x,y\in K$ with $d(x,y)\leq r_k$.
They join inside $B(x,Lr_k)$. Choose $b$ with $a^b>2L$ and consider
indices $l\leq k-b$ above a fixed interior threshold, in an arithmetic
progression of step $b_1$ with $a^{b_1}>16$. The annuli
$\{r_l/4\leq d(x,\cdot)\leq4r_l\}$ are disjoint. For $N$
sufficiently large relative to $k$, \eqref{eq:annular} and
\eqref{eq:records} give
\[
 |H_N(x)-H_N(y)|^2\leq C\lambda^{-l}
 \mu_N\{r_l/4\leq d(x,\cdot)\leq4r_l\}.
\]
Summing over at least $ck$ indices yields
$|H_N(x)-H_N(y)|\leq C_Kk^{-1/2}$. If $\lambda>1$, the single
choice $l=k-b$ gives $C_K\lambda^{-k/2}$ instead.

For each tolerance this controls a tail uniformly; the finitely many
omitted functions are individually uniformly continuous. Thus the
sequence is equicontinuous on $K$. Arzel\`a--Ascoli on a compact
exhaustion of $U$, followed by weak compactness of the measures on
$Z$, gives the asserted subsequence. The two constant band values
prove nonconstancy. The stronger estimate gives local H\"older exponent
$\log\lambda/(2\log a)$.

For \eqref{eq:measure}, fix $A>2L$ and choose
$As\leq R=r_l<aAs$. Any two points of $\overline B(x,s)$ join in
$B(x,Ls)\subset B(x,R/2)$. The annular and record estimates give
\[
 \bigl(\osc_{\overline B(x,s)}H_N\bigr)^2
       \leq C\mu_N\bigl(\overline B(x,4aAs)\bigr)
\]
for large $N$. Local uniform convergence passes the oscillation to
the limit; the closed-set inequality for weak convergence bounds the
right-hand limsup by $C\mu(\overline B(x,4aAs))$.
Take $C_\mu=4aA$.

For each fixed $t$, take a Hausdorff limit $C_t$ of $C_{N,t}$ in
$\cQ$. It is connected and has diameter at least $d_1$. Local
uniform convergence gives $u=t$ on $C_t\cap U$. A fixed compact
straddling path also straddles $t$ for every sufficiently large
$H_N$, hence meets those $C_{N,t}$ and their limit. The subsequence
used here may depend on $t$; no simultaneous selection is needed.
\end{proof}

\subsection{Two comparison tools}

The scalar construction will be used in three ways. The large continua
$C_t$ provide levels for annular constructions, with $u=t$ on
$C_t\cap U$; no value of $u$ is asserted outside $U$. The measure
estimate \eqref{eq:measure} turns oscillation on a small ball into a
lower bound for its enlarged ball's $\mu$-mass. To compare
transformations, we retain $H_N$ and $C_{N,t}$: the identity $H_N=t$
holds on the entire continuum, and the local oscillation energies
tend to zero after each transformation has been fixed. Those
comparisons will pass to subcontinua contained in compact subsets of
the common open domains.

The first tool integrates variation of a second function over a
family of level continua. It is a finite-cell statement, independent
of the preceding construction.

\begin{lemma}[A finite-cell level comparison]
\label{lem:cross-energy}
Let $f,w$ be continuous on an open set $V\subset Z$, let $K\subset V$
be compact, and let $I\subset\RR$ be an interval. Suppose that for
every $t\in I$ there is a continuum $D_t\subset K$ with $f=t$ on
$D_t$ and $\osc_{D_t}w\geq b_0>0$. For sufficiently large $n$,
\begin{equation}
 b_0|I|\leq
 \sum_{\substack{q\in T_n\\B(q,2r_n)\cap K\ne\varnothing}}
 \osc_{\overline B(q,2r_n)}f\,
 \osc_{\overline B(q,2r_n)}w.
 \label{eq:cross-energy}
\end{equation}
Consequently these hypotheses cannot hold for sequences with both
local oscillation energies on $K$ tending to zero, if $b_0$ and
$|I|$ stay bounded below by positive constants.
\end{lemma}
\begin{proof}
Let $I_q(f)$ be the closed range interval of $f$ on the closed cell.
For large $n$ these cells lie in $V$. The interval $w(D_t)$ is covered
by the range intervals of $w$ on cells meeting $D_t$, and each such
cell has $t\in I_q(f)$. Therefore
\[
 b_0\leq\sum_{q:B(q,2r_n)\cap K\ne\varnothing}
 \osc_{\overline B(q,2r_n)}w\,\mathbf1_{I_q(f)}(t).
\]
Integrate over $I$ and use
$|I\cap I_q(f)|\leq\osc_{\overline B(q,2r_n)}f$.
Only finitely many interval indicators are integrated, so no
measurable selection of $D_t$ is required. Cauchy--Schwarz gives
the final assertion.
\end{proof}

\begin{lemma}[Stopping a continuum at a sphere]
\label{lem:stop-continuum}
Let $D$ be a continuum, $p\in Z$, and $r>0$. If
$x\in D\setminus\overline B(p,r)$ and $D\cap B(p,r)\ne\varnothing$,
then $D$ contains a subcontinuum through $x$ meeting the radius-$r$
sphere and avoiding $B(p,r)$.
\end{lemma}
\begin{proof}
In a compact metric space $X$, a point's component is the intersection
$Q_x$ of its clopen neighborhoods. To verify connectedness of $Q_x$,
suppose it is separated into two compact pieces in disjoint open sets
$O_1,O_2$. Compactness of $X\setminus(O_1\cup O_2)$ gives a finite
intersection $W$ of clopen neighborhoods of $x$ contained in their
union. Its intersection with the $O_i$ containing $x$ is clopen in
$X$ and misses the other piece of $Q_x$, a contradiction. The same
compactness argument shows that every open neighborhood of a component
contains a clopen neighborhood of it.

Apply this to $X=D\setminus B(p,r)$ and the component $C$ of $x$.
If $C$ missed the sphere, it would have a clopen neighborhood $W$ in
$X$ missing the sphere. Compactness places $W$ a positive distance
from $\overline B(p,r)$, so $W$ is also clopen in $D$. It is nonempty
and misses $D\cap B(p,r)$, contradicting connectedness. Thus $C$ is
the required subcontinuum.
\end{proof}

\section{Exponential records: comparing two magnifications}
\label{sec:exponential}

Suppose $\lambda>1$. The scalar limit $u$ is then locally H\"older.
We magnify its variation to a nonconstant function on a punctured
sphere. The approximating level continua force every transform with a different
puncture to be constant wherever this function varies. A second
magnification makes two such transforms vary on the same two disks,
with disjoint value ranges, yielding a contradiction.
Scalar approximation indices throughout belong to the subsequence
of \Cref{prop:scalar-limit}.

\subsection{Magnifying variation}

\begin{lemma}[H\"older magnification]
\label{lem:holder-magnification}
Let $f$ be a real-valued, locally H\"older continuous function on an open set $\Omega\subset Z$.
Suppose $O$ is open, $\overline O\subset\Omega$, and $O$ contains a
compact path on which $f$ is nonconstant. Fix $z_0$ with
$\dist(z_0,O)>0$. There are depths $k_i\to\infty$, centers $x_i\in O$,
positive numbers $\sigma_i$, and expansions $h_i=v_i^{-1}$ at
$(x_i,k_i)$ such that, after passage to a subsequence,
\[
 p_i=h_i z_0\longrightarrow b,
 \qquad
 \frac{\dist(x_i,Z\setminus O)}{r_{k_i}}\longrightarrow\infty.
\]
Every compact subset of $Z\setminus\{b\}$ eventually lies in
$h_i\Omega$, and
\begin{equation}
 F_i(z)=\frac{f(h_i^{-1}z)-f(x_i)}{\sigma_i}
 \label{eq:blowup}
\end{equation}
converges locally uniformly there to a nonconstant locally H\"older
continuous function $F$.
\end{lemma}

\begin{proof}
Write $\delta(x)=\dist(x,Z\setminus O)$ and fix $\eta>Q$.
A finite cover of $\overline O$ gives a uniform H\"older bound,
with some exponent $\alpha_0>0$, on a neighborhood of $\overline O$.
For all large $k$ put
\[
 A_k=\sup_{x\in O}\delta(x)^\eta
            \osc_{\overline B(x,r_k)}f.
\]
The balls lie in $\Omega$, and
\begin{equation}
 c a^{-Qk}\le A_k\le C a^{-\alpha_0 k}.
 \label{eq:magnification-bounds}
\end{equation}
For the lower bound, cover a compact variation path $P\subset O$
by at most $C r_k^{-Q}$ radius-$r_k$ balls centered on $P$, using
\Cref{lem:packing}. Their range intervals cover the interval $f(P)$,
so one ball has oscillation at least $c r_k^Q$; also
$\min_P\delta>0$. The upper bound is the H\"older estimate.

Fix $0<\alpha<\alpha_0$. Since $A_k a^{\alpha k}\to0$, there are
arbitrarily late indices attaining their maximum over the remaining
tail. At these indices,
\begin{equation}
 A_{k+j}\le a^{-\alpha j}A_k\qquad(j\ge0).
 \label{eq:magnification-tail}
\end{equation}
Choose $x_k\in O$ attaining at least half the supremum, and set
$\delta_k=\delta(x_k)$ and $\sigma_k=A_k/\delta_k^\eta$.
Boundedness of $f$ and \eqref{eq:magnification-bounds} give
$\delta_k^\eta\ge c r_k^Q$, hence
$\delta_k/r_k\ge c r_k^{Q/\eta-1}\to\infty$.
Expand at $(x_k,k)$ and take a subsequence with $p_k=h_kz_0\to b$.

For a compact $K\subset Z\setminus\{b\}$, \eqref{eq:puncture} gives
\begin{equation}
 h_k^{-1}K\subset\overline B(x_k,C_Kr_k)\subset O
 \label{eq:magnification-preimage}
\end{equation}
eventually. The expansion formula gives
$d(h_k^{-1}z,h_k^{-1}w)\le C_Kr_kd(z,w)$ on $K$.
On every fixed multiple of this source ball, $\delta$ is comparable
to $\delta_k$. Thus \eqref{eq:magnification-tail} implies
\begin{equation}
 \frac{|f(y)-f(y')|}{\sigma_k}\le C a^{-\alpha j}
 \quad\text{if }d(y,y')\le r_{k+j},\quad j\ge0.
 \label{eq:magnification-local-holder}
\end{equation}
For sufficiently small $d(z,w)$, choosing $j\ge0$ with
$a^{-j}\ge C_Kd(z,w)>a^{-j-1}$ proves a uniform local
$\alpha$-H\"older estimate for $F_k$ on $K$.

For boundedness, \Cref{lem:paths} joins $h_k^{-1}z$ to $x_k$
inside a fixed multiple of $B(x_k,r_k)$. Packing covers this path
by a bounded number of radius-$r_k$ balls centered on the path, each with oscillation
at most $C\sigma_k$. Their range intervals cover the path's image,
so $|F_k(z)|\le C_K$. Compact exhaustion and Arzel\`a--Ascoli now
give a locally uniform, locally H\"older limit.
Finally, $\osc_{h_k\overline B(x_k,r_k)}F_k\ge1/2$.
The lower puncture bound places all these image balls in one compact
set off $b$, so local uniform convergence makes the limit nonconstant.
\end{proof}

\subsection{A common mesh and escaping level continua}

Apply \Cref{lem:holder-magnification} to $u$ on $U$, taking a small
variation region $O$ and $z_0$ outside its closure. Call the limit
$W$, its puncture $b$, and retain $x_i,k_i,h_i,\sigma_i$.

\begin{lemma}[Diagonal approximants]
\label{lem:exponential-diagonal}
After thinning the expansions, there are increasing scalar indices
$N_j$ such that
\[
 f_j(z)=\frac{H_{N_j}(h_j^{-1}z)-u(x_j)}{\sigma_j},
 \qquad z\in h_j\Qhat,
\]
converges locally uniformly to $W$ on $Z\setminus\{b\}$.
For every fixed $g\in G$, compact $K\subset Z\setminus\{gb\}$,
and fixed $J\ge2$,
\begin{equation}
 \mathcal E_{N_j}^{(J)}(f_j\circ g^{-1};K)\longrightarrow0.
 \label{eq:exponential-translated-energy}
\end{equation}
All indicated enlarged balls eventually lie in the functions' domains.

If a compact path $P\subset Z\setminus\{b\}$ has endpoint values
of $W$ strictly straddling a compact positive-length interval $I$,
then, for all large $j$ and every $t\in I$, there is a continuum
$D_{j,t}\subset h_j\Qhat$ with
\[
 f_j=t\text{ on }D_{j,t},\qquad D_{j,t}\cap P\ne\varnothing,
 \qquad \sup_{t\in I}\dist(b,D_{j,t})\longrightarrow0.
\]
\end{lemma}

\begin{proof}
For $j\ge4$ put
$K_j=\{d(z,b)\ge1/j\}$ and $L_j=\{d(z,b)\ge1/(2j)\}$.
Choose the $j$th expansion so that $h_j^{-1}L_j\subset U$ and its
normalized $u$-transform differs from $W$ by at most $1/j$ on $L_j$.
\emph{Hold this expansion and its normalization fixed.} Then choose
$N_j>N_{j-1}$ so large that
\[
 \sup_{L_j}\frac{|H_{N_j}(h_j^{-1}z)-u(h_j^{-1}z)|}{\sigma_j}
 \le\frac1j,
 \qquad \mathcal E_{N_j}^{(j)}(f_j;K_j)\le\frac1j,
 \qquad (j+2)r_{N_j}<\frac1{2j}.
\]
The uniform approximation follows from \Cref{prop:scalar-limit}.
With $h_j$ and $\sigma_j$ fixed, \Cref{lem:local-oscillation} bounds
the displayed energy by $C(h_j,K_j,j)/(\sigma_j^2M(N_j))$, so the
energy condition follows by increasing $N_j$. The last condition
keeps all the enlarged balls in $L_j$. This proves local uniform
convergence and vanishing energy at every fixed enlargement off $b$.

For a fixed $g$, choose a compact neighborhood $K'$ of $g^{-1}K$
avoiding $b$. Pulling back each target ball and choosing a nearest
source net point gives, by bi-Lipschitz distortion and packing,
\[
 \mathcal E_n^{(J)}(F\circ g^{-1};K)
 \le C_g\mathcal E_n^{(J')}(F;K')
\]
for some fixed $J'\ge2$ and all large $n$. Each source center receives
only boundedly many target centers. Eventually $K'\subset K_j$ and
$J'\le j$, proving \eqref{eq:exponential-translated-energy}.

Uniform convergence makes the endpoint values of $f_j$ straddle $I$.
The source level $\tau_j(t)=u(x_j)+\sigma_jt$ therefore lies in $(0,1)$,
and \Cref{lem:level-continua} supplies
$D_{j,t}=h_jC_{N_j,\tau_j(t)}$, meeting $P$.
Its source continuum has diameter at least $d_1$, so contains a point
at distance at least $d_1/2$ from $x_j$. By \eqref{eq:puncture},
\[
 \dist(b,D_{j,t})\le d(b,p_j)+Cr_{k_j}\longrightarrow0
\]
uniformly for $t\in I$.
\end{proof}

\subsection{Comparison across punctures}

Let $\mathcal N(F)$ be the points of a continuous function's domain
at which it is not locally constant.

\Needspace{5\baselineskip}
\begin{lemma}[Transform comparison]
\label{lem:exponential-comparison}
If $c=gb\ne b$ and $V=W\circ g^{-1}$, then
\begin{equation}
 V(x)=V(b)\qquad
 \bigl(x\in\mathcal N(W)\setminus\{c\}\bigr).
 \label{eq:transform-level}
\end{equation}
More generally, for $W_s=W\circ s^{-1}$ and $p_s=sb$, if $p_s\ne p_t$,
\begin{equation}
 W_t(x)=W_t(p_s)\qquad
 \bigl(x\in\mathcal N(W_s)\setminus\{p_t\}\bigr).
 \label{eq:two-transform-level}
\end{equation}
\end{lemma}

\begin{proof}
Suppose $x\in\mathcal N(W)\setminus\{c\}$ and $V(x)\ne V(b)$.
Choose $\varepsilon>0$ and a coordinate disk about $x$, compactly
contained off $b,c$, on which $|V-V(b)|\ge4\varepsilon$.
It contains a compact variation path $P$ for $W$. Choose a compact
positive-length interval $I$ strictly between its endpoint values,
avoiding $W(c)$. Continuity gives $\epsilon_c>0$ such that
$\overline B(c,\epsilon_c)$ avoids $b$ and its $W$-range has positive
distance from $I$. Consequently, for all large $j$, the continua
$D_{j,t}$ from \Cref{lem:exponential-diagonal} avoid this closed ball.

Choose $\rho>0$ so that $\overline B(b,\rho)$ avoids $c,P$ and
$|V-V(b)|\le\varepsilon$ there. For $w_j=f_j\circ g^{-1}$,
local uniform convergence gives
\[
 |w_j-V(b)|\ge3\varepsilon\quad\text{on }P,
 \qquad |w_j-V(b)|\le2\varepsilon\quad\text{on }\overline B(b,\rho)
\]
for all large $j$. Each $D_{j,t}$ meets $P$ and enters $B(b,\rho)$.
By \Cref{lem:stop-continuum}, it has a subcontinuum joining a chosen
point of $P$ to the radius-$\rho$ sphere and contained in
\[
 K=\{z:d(z,b)\ge\rho,\ d(z,c)\ge\epsilon_c\}.
\]
This subcontinuum lies in the level $f_j=t$ and has $w_j$-oscillation
at least $\varepsilon$. A fixed compact neighborhood of $K$ avoiding
both punctures eventually lies in $h_jU\cap gh_jU$ and contains every
closed radius-$2r_{N_j}$ cell meeting $K$. Hence
\Cref{lem:cross-energy} and Cauchy--Schwarz give
\[
 \varepsilon|I|\le
 \bigl(\mathcal E_{N_j}^{(2)}(f_j;K)\bigr)^{1/2}
 \bigl(\mathcal E_{N_j}^{(2)}(w_j;K)\bigr)^{1/2}\longrightarrow0,
\]
a contradiction. All sets and $g$ are fixed in this limit.
Finally apply \eqref{eq:transform-level} at $s^{-1}x$ with group
element $s^{-1}t$ to obtain \eqref{eq:two-transform-level}.
\end{proof}

\subsection{The second magnification}

We first control repeated punctures. An accumulation direction of
vertices going to infinity is a subsequential boundary limit.

\begin{lemma}[Stabilizer directions]
\label{lem:stabilizer}
For every $b\in Z$, the vertices of $E=\Stab_G(b)$ have at most
two accumulation directions.
\end{lemma}

\begin{proof}
Let $\xi$ be a based ray to $b$. For $q\in E$, $|q|\le t$,
thinness of the ideal triangle with sides $[e,q],\xi,q\xi$ puts
$q\xi(n)$ uniformly near $\xi$ when $n$ is sufficiently larger
than $t$. The nearby parameter lies in $[n-t-C,n+t+C]$.
For one common integer $n$, these vertices are distinct as $q$ varies.
Bounded valence therefore gives
\begin{equation}
 \#\{q\in E:|q|\le t\}\le C'(t+1).
 \label{eq:stabilizer-growth}
\end{equation}

If another accumulation direction exists, take $q_n\to x\ne b$ and
$q_n^{-1}\to y$ along a subsequence. Base change gives, uniformly
on compact sets avoiding $y$,
\[
 (q_n|q_nz)_e=|q_n|-(q_n^{-1}|z)_e+O(1)\ge |q_n|-O(1).
\]
By \Cref{lem:prefix}, $q_nz\to x$ there, and $q_nb=b$ forces $y=b$.
Choose disjoint small neighborhoods of $b,x$, with a point outside
their union. For large $n$, $q=q_n$ sends the complement of the
first into the second. Its positive iterates never return the outside
point, so $q$ has infinite order. It is therefore loxodromic, with
endpoints $b,x'$; see \cite[III.$\Gamma$.3.10]{BH}.

If some $h\in E$ moved $x'$, the cyclic subgroups generated by $q$
and $hqh^{-1}$ would intersect trivially: a nonzero common power
would have the same pair of endpoints. Their $m^2$ products
$q^i(hqh^{-1})^j$, $1\le i,j\le m$, would then be distinct and have
length $O(m)$, contradicting \eqref{eq:stabilizer-growth}.
Thus every element of $E$ fixes both endpoints. It maps a geodesic
line between them to a line at uniformly bounded Hausdorff distance,
so its orbit of $e$ stays a fixed distance from that line. Its only
possible accumulation directions are the two endpoints.
\end{proof}

\begin{proposition}
\label{prop:no-exponential-records}
The contradiction setup cannot have $\lambda>1$.
\end{proposition}

\begin{proof}
Let $E=\Stab(b)$ and let $\Lambda_E$ be its at-most-two accumulation
directions. Choose a variation region $O$ for $W$ whose closure avoids
$\{b\}\cup\Lambda_E$. Such a region exists: otherwise $W$ would
be locally constant on the connected sphere minus this finite set,
and continuity would make it constant throughout its domain.

Apply \Cref{lem:holder-magnification} to $W$ on $Z\setminus\{b\}$,
with this $O$ and $z_0=b$. Write the expansions as $s_i=v_i^{-1}$,
with centers $y_i\in O$, depths $\ell_i\to\infty$, and punctures
$p_i=s_i b\to b_*$. Positive affine normalizations of
$W_i=W\circ s_i^{-1}$ converge locally uniformly off $b_*$ to a
continuous nonconstant $W_*$.

The punctures have a pairwise distinct subsequence. Indeed, if
$v_i^{-1}b=v^{-1}b$ infinitely often, then $v_i v^{-1}\in E$.
After a subsequence $y_i\to y\in\overline O$. Since $v_i$ lies at
depth $\ell_i$ on a ray to $y_i$, \Cref{lem:prefix} gives $v_i\to y$
in boundary direction. Bounded right multiplication preserves this
direction, implying $y\in\Lambda_E$, a contradiction.

Choose two coordinate disks $D_1,D_2$, with compact closures off $b_*$,
having disjoint $W_*$-range intervals and each containing a compact
variation path. To do this, follow a variation path to its first
hitting points of two distinct intermediate values. Just before either
first hitting time the values differ, so the function is not locally
constant at either point. Continuity supplies disks with disjoint
ranges, and non-local constancy supplies a variation path in each.

For every sufficiently large $i$, the puncture $p_i$ is outside the
closed disks. Uniform convergence of the normalized $W_i$ preserves
both the disjoint range intervals and the positive variation on the
two fixed paths. Positive affine normalization preserves these
properties, so $\mathcal N(W_i)$ meets each disk.
Fix two distinct such indices $i,j$. By
\eqref{eq:two-transform-level}, $W_j$ has the single value $W_j(p_i)$
at every point of $\mathcal N(W_i)$ in either disk. Taking one point
in each contradicts the disjoint $W_j$-range intervals.
\end{proof}

\section{A probability flow in the subexponential case}
\label{sec:flow}

Assume now that the unbounded moduli satisfy
\begin{equation}
 \limsup_{j\to\infty}M(j)^{1/j}\leq1.
 \label{eq:subexponential-growth}
\end{equation}
Let $u,\mu,C_t$ be given by \Cref{prop:scalar-limit}. We build a tree of
small balls carrying nested intervals of values of $u$. At each node,
either two separated variations survive expansion, or many children
carry a substantial fraction of its interval. A probability flow will
force the first alternative at arbitrarily deep levels.

Choose $x_0\in U$ where $u$ is not locally constant, and choose
$\epsilon_0>0$ with $\overline B(x_0,3\epsilon_0)\subset U$. Put
\[
 K=\overline B(x_0,\epsilon_0),\qquad
 K^+=\overline B(x_0,2\epsilon_0),
\]
and, shrinking $\epsilon_0$ if necessary, choose $z_0$ a positive distance
from $K^+$. All centers will lie in $K$ and all source neighborhoods in
$K^+$. Use the constants of \eqref{eq:measure} for this fixed region.

\subsection{Node intervals and transition probabilities}

The tree will record both where variation occurs and which values it
carries. A node has a source center $x$, a scale $R$, and a value interval
$J$. Its children have scale $r=\theta R$, where $\theta=a^{-D}$ for
one positive integer subdivision depth $D$ used throughout the tree. Their intervals
are contained in $J$. We will choose probabilities on the children
that control interval length linearly and quadratically; the two
controls have different roles in the final comparison.

We first define this data for geometric constants $c_0>2L+2$,
$Lc_0+2<A<B_0-2$, and $H>0$. Here $c_0$ is the source-ball enlargement,
$A,B_0$ locate an annulus of possible child centers, and $Hr$ is the
separation required of retained children. The next subsection fixes
these constants, before choosing $D$. We use only scales $R=r_k$
small enough that all fixed enlargements below lie in $K^+$ and
satisfy \eqref{eq:measure}, and that $(B_0+2)R<d_1/2$.

A \emph{node} is a center $x\in K$, a scale $R=r_k$, and a nonempty
open interval $J=(j_-,j_+)\subset(0,1)$ such that
\begin{equation}
 J\subset\left(\min_{\overline B(x,c_0R)}u,
                   \max_{\overline B(x,c_0R)}u\right).
 \label{eq:node-straddling}
\end{equation}
Write $\Delta=|J|$. For every $t\in J$, a path in $B(x,Lc_0R)$
strictly straddles $t$ and hence meets $C_t$. Since
$\diam C_t\geq d_1$, the continuum leaves
$\overline B(x,(B_0+2)R)$ and meets every intervening radial sphere.

Take a maximal $r$-separated set $Y$ in
\[
 \mathcal A_x(R)=\{y:(A-1)R\leq d(x,y)\leq(B_0+1)R\}.
\]
It is an $r$-net there, and packing gives
\begin{equation}
 \#Y\leq C_Ya^{QD},
 \label{eq:candidate-count}
\end{equation}
with $C_Y$ independent of $D$. For $y\in Y$ let $I_y^{\mathrm{in}}$
and $I_y^{\mathrm{out}}$ be the closed intervals between the extrema
of $u$ on $\overline B(y,2r)$ and $\overline B(y,c_0r)$, respectively.
These are range \emph{intervals}; the balls need not be connected. Put
\[
 J_y=J\cap\operatorname{int}I_y^{\mathrm{out}},\qquad
 t_y=|J_y|/\Delta.
\]
A positive $t_y$ gives a nested open child interval satisfying
\eqref{eq:node-straddling}, and, including singleton ranges,
\begin{equation}
 |J\cap I_y^{\mathrm{out}}|=t_y\Delta.
 \label{eq:closed-range-length}
\end{equation}

For $0<c_*\leq1/(8\max\{C_Y,1\})$, put $\kappa=c_*a^{-QD}$.
A node is \emph{good} if two candidates $y^-,y^+$ satisfy
\begin{align}
 |I_{y^-}^{\mathrm{in}}\cap(j_-,j_-+\Delta/4)|&\geq\kappa\Delta,
 &|I_{y^+}^{\mathrm{in}}\cap(j_-+3\Delta/4,j_+)|&\geq\kappa\Delta,
 \label{eq:good-inner}\\
 \dist(I_{y^-}^{\mathrm{out}},I_{y^+}^{\mathrm{out}})&\geq\Delta/4.
 \label{eq:good-outer}
\end{align}
Goodness uses the full candidate set, independently of the children
later retained.

The required probabilities are as follows. Fix the exponents
\begin{equation}
 b_* =\frac12,\qquad \tau=\frac1{20},\qquad
 \beta=\frac{\tau}{4(Q+1)},\qquad \gamma=\frac\beta2.
 \label{eq:flow-exponents}
\end{equation}
At every node we seek an $Hr$-separated set of children, with
probabilities summing to one, such that
\begin{align}
 p_y&\leq a^{-\gamma D}t_y,
 &p_y&\leq a^{-b_*D}t_y^2
 &&\text{at nodes that are not good},\label{eq:transition-bad}\\
 p_y&\leq a^{-\gamma D}t_y,
 &p_y&\leq a^{\tau D}t_y^2
 &&\text{at good nodes}.\label{eq:transition-general}
\end{align}
Multiplying the linear bounds along a branch will telescope the
interval ratios and make the branch probability small. Multiplying
the quadratic bounds will compare probability to squared interval
length. The measure estimate bounds the sum of those squares at
each separated generation. It will therefore bound the average
probability of nodes with the stronger quadratic bound and force a
substantial amount of good-node probability at arbitrarily deep levels.

Two facts will supply these transitions. A node that is not good has
many children with $t_y\ge1/8$, on which uniform probabilities suffice.
At every node, a capacity estimate supplies enough total weight to
satisfy the weaker pair of bounds \eqref{eq:transition-general}.
After constructing the transitions, we show that goodness also gives
two separated regions of variation in one of finitely many target
patterns. That geometric property is the input to \Cref{sec:pairs}.

\subsection{Geometric choices and many large children}

We now fix the geometry used by the definitions above. Put
$c_{\mathrm{in}}=2L$ and choose $c_0>c_{\mathrm{in}}+2$ so large that,
with
\[
 A=4Lc_0+20,\qquad B_0=A+10,\qquad
 c_{\mathrm{pat}}=\frac{2C_Ec_{\mathrm{in}}}{(A-2)^2},
\]
we have
\begin{equation}
 \frac{c_Ec_0}{(B_0+2)^2}>12c_{\mathrm{pat}}.
 \label{eq:flow-constant-choice}
\end{equation}
The constant $c_{\mathrm{pat}}$ will be the target radius factor:
paths joining extrema of the two inner child balls will expand into balls of radius
$c_{\mathrm{pat}}\theta/2$, while the larger source balls will contain
the preimages of radius-$10c_{\mathrm{pat}}\theta$ target balls.
The ratio of the left side of \eqref{eq:flow-constant-choice} to the
right side grows linearly in $c_0$, so this choice is possible.
Next choose
\[
 H>\max\{4(B_0+1),4C_{\mu}c_0,2\}.
\]
The first inequality preserves separation between different branches;
the second makes the measure enlargements at one generation disjoint.
These choices fix $C_Y$ in \eqref{eq:candidate-count}, and we fix
$c_*$ as specified above. Require $\theta<1/4$ and $c_0\theta<1$.
All further lower bounds on $D$ will depend only on this fixed data.

The uniform upper bound on node scales can be chosen now, before $D$:
all source radius factors are fixed, and the subsidiary scales $r$
and $s$ used below never exceed $R$. Thus the source neighborhoods
and measure enlargements stay in $K^+$ for every node. We choose the
root scale after $D$ has been fixed.

\begin{lemma}[Many children at a node that is not good]
\label{lem:bad-node-children}
For sufficiently large $D$, every node that is not good has at least
$c_Ha^D$ mutually $Hr$-separated candidates with $t_y\geq1/8$,
where $c_H>0$ is independent of the node and $D$.
\end{lemma}
\begin{proof}
Fix $b\in[AR,B_0R]$. Each $C_t$, $t\in J$, meets the radius-$b$
sphere, so the inner intervals of candidates with
$|d(x,y)-b|\leq r$ cover $J$. By \eqref{eq:candidate-count}, two of
them meet the respective quarters in \eqref{eq:good-inner} by at least
$\kappa\Delta$.

If both have $t_y<1/8$, the lower-quarter outer interval ends below
$j_-+3\Delta/8$, and the upper-quarter interval starts above
$j_-+5\Delta/8$. Indeed each meets its quarter but overlaps $J$ in
length less than $\Delta/8$; the assertion also holds for intervals
extending beyond $J$. Their distance exceeds $\Delta/4$, making the
node good. Otherwise one candidate has $t_y\geq1/8$. Apply this at
radial values spaced by $(H+3)r$: there are at least $c_Ha^D$ such
values, and the resulting centers are more than $Hr$ apart.
\end{proof}

At a node that is not good, put uniform probabilities on the
candidates in \Cref{lem:bad-node-children}. For large $D$ they satisfy
\eqref{eq:transition-bad}: indeed, $p_y\le c_H^{-1}a^{-D}$ and
$t_y\ge1/8$, while $\gamma,b_*<1$.

\subsection{The capacity estimate}

It remains to construct transitions satisfying
\eqref{eq:transition-general}. A single subdivision depth $D$ must
work even when a node's interval $J$ is arbitrarily short. We therefore
integrate over $J$ before choosing $D$; its length will cancel from
the estimate.

\begin{lemma}[Separated capacity]
\label{lem:separated-capacity}
For sufficiently large $D$, uniformly over all nodes, there is an
$Hr$-separated subset $Y'\subset Y$ such that
\begin{equation}
 \sum_{y\in Y'}\min\{a^{-\gamma D}t_y,a^{\tau D}t_y^2\}\geq1.
 \label{eq:capacity}
\end{equation}
\end{lemma}
\begin{proof}
Join two candidates in a graph when their distance is at most $Hr$.
Packing bounds its degree, so graph coloring partitions $Y$ into at most
$C_H$ separated classes, with $C_H$ independent of $D$. If the conclusion fails, then
\begin{equation}
 \sum_{y\in Y}\min\{a^{-\gamma D}t_y,a^{\tau D}t_y^2\}<C_H.
 \label{eq:failed-capacity}
\end{equation}
Call $y$ \emph{large} when $t_y\geq a^{-(\tau+\gamma)D}$ and
\emph{small} otherwise. On large candidates the minimum is linear, so
\begin{equation}
 \sum_{\text{large }y}t_y\leq C_Ha^{\gamma D}.
 \label{eq:large-total-length}
\end{equation}

We will test levels at an intermediate radius $s$, larger than the
child radius $r$ and smaller than the parent radius $R$. This provides
more radial tests than the large ranges can occupy. At the same time,
it lets us detect the remaining levels with a modulus cost at the
relative scale $r/s$.

\emph{Remove the large ranges.}
Set $D'=\lceil\beta D\rceil$ and $s=a^{-D'}R$. For large $D$,
$0<D'<D$ and $s<R/4$. Choose $s$-spaced radial bins $\mathcal B$
in $[(A+1)R,(B_0-1)R]$, so $\#\mathcal B\asymp a^{D'}$.
On $\mathcal B\times J$ use counting measure times Lebesgue measure.
Mark $(b,t)$ if a large candidate satisfies
\[
 |d(x,y)-b|\leq5s,\qquad t\in I_y^{\mathrm{out}}.
\]
Each candidate meets at most 11 bins. By
\eqref{eq:closed-range-length} and \eqref{eq:large-total-length}, the
marked measure is at most $CC_Ha^{\gamma D}\Delta$. Since $\gamma<\beta$,
the unmarked measure is at least
\begin{equation}
 ca^{D'}\Delta.
 \label{eq:unmarked-mass}
\end{equation}

\emph{Detect the remaining levels by modulus.}
Take a maximal $s$-separated set $\mathcal Z_s$ in $\mathcal A_x(R)$;
its radius-$s$ balls cover that annulus, and it has at most $Ca^{QD'}$ points. Expand at
each $z\in\mathcal Z_s$ at depth $k+D'$ by $h_z$. Put
$q=r/s=a^{-(D-D')}$, and take macroscopic admissible weights
$\omega_\xi$, $\xi\in T_{D-D'}$, of squared mass $M(D-D')$.
For fixed $C_T>2+C_E$, define
\[
 \sigma_y^z=
 \begin{cases}
 \displaystyle\sum_{\xi:\ d(\xi,h_zy)\leq C_Tq}\omega_\xi,
       &d(y,z)\leq3s,\\[3pt]
 0,&d(y,z)>3s,
 \end{cases}
 \qquad \sigma_y=\sum_{z\in\mathcal Z_s}\sigma_y^z.
\]
Within $\overline B(z,3s)$ the candidate images are separated by a
fixed multiple of $q$. Thus every sum has boundedly many terms and
every target weight is charged boundedly many times. Thus
$\sum_y(\sigma_y^z)^2\leq CM(D-D')$ for each $z$. Also each $y$
lies within $3s$ of boundedly many auxiliary centers. Packing and
Cauchy--Schwarz give
\begin{equation}
 \sum_y\sigma_y^2\leq Ca^{QD'}M(D-D').
 \label{eq:capacity-transfer-energy}
\end{equation}

Fix an unmarked $(b,t)$. Choose $w\in C_t$ on the radius-$b$ sphere
and $z\in\mathcal Z_s$ with $d(w,z)\leq s$. The diameter bound gives
a point of $C_t$ outside $\overline B(w,2s)$. Join these points in the
$r/4$-neighborhood of $C_t$, by \Cref{lem:paths}, and stop at the first
radius-$s$ hit about $w$. The resulting path $\pi$ has diameter at
least $s$ and lies in $\overline B(w,s)\subset\overline B(z,2s)$.
The scale-$s$ expansion $h_z$ therefore makes it macroscopic by the
choice of $d_0$ in \Cref{sec:geometry}.

The radial bin lies between $(A+1)R$ and $(B_0-1)R$, while $s<R/4$.
Thus every point of $\pi$ lies between radii $(A+3/4)R$ and
$(B_0-3/4)R$ about $x$. In particular it lies in the candidate
annulus and has a candidate center within distance $r$. The same
fixed source margin places every level point within $r/4$ of $\pi$
in $K^+\subset U$. On these points the scalar-limit property gives
$u=t$, even though it does not assert this identity on parts of
$C_t$ outside $U$.

For a target cell meeting $h_z\pi$, choose a preimage hit $v\in\pi$
and $y\in Y$ with $d(y,v)\leq r$. Then $d(y,z)<3s$. If $\xi$ is the target cell's center,
\eqref{eq:expansion} gives $d(\xi,h_zy)<(2+C_E)q<C_Tq$, so its
weight occurs in $\sigma_y^z$. A point of $C_t$ is within $r/4$ of
$v$, hence within $5r/4<c_0r$ of $y$. The preceding interior check
gives $t\in I_y^{\mathrm{out}}$. Since
$|d(x,y)-b|\leq s+r<5s$, unmarkedness makes $y$ small. Admissibility
therefore yields
\begin{equation}
 \sum_{\substack{\text{small }y:\ |d(x,y)-b|\leq5s}}
 \sigma_y\mathbf1_{I_y^{\mathrm{out}}}(t)\geq1
 \quad\text{on every unmarked pair}.
 \label{eq:unmarked-indicator}
\end{equation}

The unmarked levels now have an admissible detection weight supported
on small candidates. We compare its squared mass with the surviving
amount of level parameter.

\emph{Compare the energies.}
Integrating \eqref{eq:unmarked-indicator}, using bounded bin multiplicity,
gives
\[
 ca^{D'}\Delta\leq C\Delta\sum_{\text{small }y}\sigma_y t_y.
\]
Only finitely many closed-interval indicators are integrated; no
measurable selection of continua or paths is needed, and singleton
ranges contribute zero. Cancel $\Delta>0$. Cauchy--Schwarz and
\eqref{eq:capacity-transfer-energy} imply
\begin{equation}
 \sum_{\text{small }y}t_y^2
 \geq\frac{ca^{-(Q-2)D'}}{M(D-D')}.
 \label{eq:small-square-capacity}
\end{equation}
Subexponential growth gives $M(D-D')\leq a^{\tau D/4}$ for large $D$: first
bound $M(j)\leq Ca^{\tau j/8}$ and then absorb $C$ by increasing $D$. Since $D'\leq\beta D+1$ and $(Q-2)\beta<\tau/4$,
\[
 \sum_{\text{small }y}a^{\tau D}t_y^2
 \geq ca^{[3\tau/4-(Q-2)\beta]D}
 \geq ca^{\tau D/2}.
\]
On small candidates this is exactly the capacity, contradicting
\eqref{eq:failed-capacity}. Every threshold for $D$ is uniform over nodes.
\end{proof}

Fix such a $D$. At each good node, divide the positive capacities on a
set $Y'$ from \Cref{lem:separated-capacity} by their sum. These
probabilities satisfy \eqref{eq:transition-general}. At the remaining
nodes use the uniform rule above. Every positive transition retains
a nonempty open child interval $J_y\subset J$.

\subsection{Common target patterns}

The probabilities are now fixed. We verify the additional geometric
information carried by a good node: its two source variations expand
into a pair of balls chosen from a finite list, with separated range
intervals. The list depends on the fixed $D$, but not on the node's
generation or interval length.

Fix a node's expansion $h=v^{-1}$ when constructing it. At a good node,
join the extrema of each inner ball by a path in
$B(y^\pm,c_{\mathrm{in}}r)$. Each path has endpoint variation at least
$\kappa\Delta$. Its points lie at distance at least $(A-2)R$ from $x$,
so \eqref{eq:expansion} puts its image within
$c_{\mathrm{pat}}\theta/2$ of $hy^\pm$.
Conversely, write $D_w=\max\{R,d(x,w)\}$. If $d(y,z)\geq c_0r$, then
$D_z\leq D_y+d(y,z)$ and monotonicity of $s/(D_y+s)$ give
\[
 d(hy,hz)\geq\frac{c_ERc_0r}{D_y(D_y+c_0r)}
 \geq\frac{c_Ec_0\theta}{(B_0+2)^2}>12c_{\mathrm{pat}}\theta.
\]
Thus
\begin{equation}
 \overline B(hy,10c_{\mathrm{pat}}\theta)
       \subset h(B(y,c_0r)).
 \label{eq:pattern-containment}
\end{equation}
Choose a fixed finite $c_{\mathrm{pat}}\theta$-net in $Z$. For each
$hy^\pm$ choose a net point within $c_{\mathrm{pat}}\theta$ and take
the closed radius-$3c_{\mathrm{pat}}\theta$ ball about it. The ordered
pair $(\mathcal B^-,\mathcal B^+)$ is the node's \emph{pattern}.
Each ball contains its path image and lies in the corresponding outer
source-ball image. Consequently, for $u\circ h^{-1}$ the two whole-ball
range intervals have distance at least $\Delta/4$, and each ball
contains a path with endpoint variation at least $\kappa\Delta$.
There are finitely many patterns, independently of depth.

\subsection{The tree and its flow}

Choose $k_0$ so large that $R_0=r_{k_0}$ satisfies all the source-scale
conditions above and
\[
 2(B_0+1)R_0<\epsilon_0/2,\qquad (B_0+2)R_0<d_1/2.
\]
Non-local-constancy at $x_0$ supplies a nonempty open interval
$J_0\subset(0,1)$ strictly inside the range on
$\overline B(x_0,c_0R_0)$. Write $\Delta_0=|J_0|$ and construct the
tree recursively, storing each node's expansion and, at good nodes,
its pattern.

A child moves by at most $(B_0+1)R$ from a parent of scale $R$.
Consequently all descendants stay within
\begin{equation}
 \frac{B_0+1}{1-\theta}R\leq2(B_0+1)R
 \label{eq:descendant-displacement}
\end{equation}
of that parent, and all centers lie in $K$. Generation $n$ has scale
$r_{k_0+Dn}$ and its centers are $H$ times that scale apart: this holds
for siblings, while children of distinct parents of scale $R$ are
separated by at least
\[
 [H-2(B_0+1)]R>HR/2>H\theta R.
\]
At any level of scale $R$, \eqref{eq:measure} bounds $\Delta_i^2$ by
$C\mu(\overline B(x_i,C_\mu c_0R))$. These enlarged balls are disjoint
because $H>4C_\mu c_0$, and hence
\begin{equation}
 \sum_i\Delta_i^2\leq E_0,\qquad E_0=C\mu(Z)<\infty.
 \label{eq:level-energy}
\end{equation}

Give the root probability 1 and let $P_i$ be the product of transitions
to node $i$. Every level has total probability 1. If $i$ has generation
$n$, let $m_i$ count its ancestors at generations $0,\ldots,n-1$ that
are not good. Multiplying the transition bounds and telescoping the
interval ratios yields
\begin{equation}
 P_i\leq a^{-\gamma Dn}\frac{\Delta_i}{\Delta_0},\qquad
 \Delta_i^2\geq\Delta_0^2P_i
 a^{D((b_*+\tau)m_i-\tau n)}.
 \label{eq:flow}
\end{equation}
Conditionally, if $i$ is $m$ generations below $j$, nesting gives
\begin{equation}
 \frac{P_i}{P_j}\leq a^{-\gamma Dm}\frac{\Delta_i}{\Delta_j}
 \leq a^{-\gamma Dm}.
 \label{eq:conditional-flow}
\end{equation}

\begin{lemma}[Levels with substantial good mass]
\label{lem:good-levels}
Infinitely many generations have total probability at least $1/2$
on good nodes.
\end{lemma}
\begin{proof}
By \eqref{eq:flow}, \eqref{eq:level-energy}, and Jensen's inequality,
\[
 \frac{E_0}{\Delta_0^2}
 \geq\sum_iP_i a^{D((b_*+\tau)m_i-\tau n)}
 \geq a^{D((b_*+\tau)\sum_iP_im_i-\tau n)}.
\]
Thus
\[
 \frac1n\sum_iP_im_i\leq\frac{\tau}{b_*+\tau}+O(1/n)
 =\frac1{11}+O(1/n).
\]
Conservation of probability identifies the left side with the average
bad mass over the first $n$ generations. If good mass were eventually
less than $1/2$, that average would have lower limit at least $1/2$,
a contradiction.
\end{proof}

\begin{proposition}[Properties of the flow]
\label{prop:probability-flow}
The construction gives a fixed $D$ and a tree with generation-$n$
centers $Hr_{k_0+Dn}$-separated and descendant displacement
\eqref{eq:descendant-displacement}. Its nonempty value intervals are
nested, its probabilities satisfy \eqref{eq:flow} and
\eqref{eq:conditional-flow}, and infinitely many levels have good mass
at least $1/2$. Every node has a fixed expansion. Each good node has
a pattern from a finite list of ordered target-ball pairs. For that expansion,
each target ball contains a path with endpoint variation at least
$\kappa\Delta_i$, their whole-ball range intervals have distance at
least $\Delta_i/4$, and the balls lie in the images of the respective
outer balls of radius $c_0\theta r_{k_0+Dn}$.
\end{proposition}

\section{Same-depth expansions and the measure comparison}
\label{sec:pairs}

We compare expansions at many nodes of a single generation. Almost
every pair of punctures is sufficiently separated. A level-continuum
argument then forces one of the two charts to spend a definite amount
of oscillation near the other's puncture. Pulling those neighborhoods
back gives more total measure than their overlap permits.

\begin{proposition}
\label{prop:no-subexponential-records}
The moduli $M(n)$ cannot be unbounded while satisfying
\[
 \limsup_{n\to\infty}M(n)^{1/n}\leq1.
\]
\end{proposition}

We prove this proposition through the pair comparison and counting
estimates below. Fix the flow of \Cref{prop:probability-flow}, including $D$, and a
compact neighborhood of its centers in $U$. Use the point $z_0$ chosen
in the flow construction, a positive distance from that neighborhood.
At generation $n$ write
\[
 k=k_0+Dn,\qquad R=a^{-k},\qquad p_i=h_i z_0,
\]
where $h_i=v_i^{-1}$ is the expansion already chosen for node $i$.
The centers $x_i$ are $HR$-separated, and
\begin{equation}
 P_i\le a^{-\gamma Dn}\frac{\Delta_i}{\Delta_0},
 \qquad \frac{P_j}{P_A}\le a^{-\gamma Dm}
 \quad\text{if $j$ lies $m$ generations below $A$}.
 \label{eq:pairs-flow-input}
\end{equation}
We use the arbitrarily large generations having good mass at least
$1/2$, supplied by \Cref{lem:good-levels}.

A good node has a pattern $(\mathcal B^-,\mathcal B^+)$ from a fixed
finite list. Each target ball contains a compact path with endpoint
$u\circ h_i^{-1}$-variation at least $\kappa\Delta_i$, and lies in
the image of its outer source ball. The two target range intervals
have distance at least $\Delta_i/4$. Throughout, a \emph{range interval}
means the interval between the minimum and maximum; metric balls need
not be connected.

All geometric and counting constants below are independent of the
generation, nodes, and approximating index $N$. Constants and thresholds
in convergence statements may depend on a pair after it has been fixed.

\subsection{A probability class and separated punctures}

\begin{lemma}[Bounded puncture multiplicity]
\label{lem:puncture-multiplicity}
For every fixed $C>0$, the balls
$\{\overline B(p_i,CR):i\text{ at generation }n\}$ have uniformly
bounded multiplicity.
\end{lemma}

\begin{proof}
Centers assigned to the same depth-$k$ vertex have mutual distance
$O(R)$. Their $HR$-separation and \Cref{lem:packing} bound their number.
If several puncture balls share a point, their punctures have mutual
products at least $k-O_C(1)$. By \eqref{eq:inverse-prefix} and
\Cref{lem:prefix}, all corresponding inverse vertices lie in one
bounded Cayley-graph ball. There are boundedly many such vertices,
each with the bounded multiplicity just proved.
\end{proof}

There are at most $Ca^{Qk}$ nodes. Discarding those with
$P_i<a^{-2Qk}$ loses probability at most $Ca^{-Qk}$, so the remaining
good nodes have mass at least $1/3$ for large $n$. Partitioning by
pattern and into the $O(k)$ dyadic probability classes gives a class
$\mathcal I$ with
\begin{equation}
 P_*\le P_i\le2P_*\quad(i\in\mathcal I),
 \qquad \sum_{i\in\mathcal I}P_i\ge\frac c k.
 \label{eq:pairs-class}
\end{equation}
Writing $\mathcal N=|\mathcal I|$, we have
\begin{equation}
 \mathcal NP_*\ge\frac{c}{2k}.
 \label{eq:pairs-class-size}
\end{equation}
Also $P_i\le a^{-\gamma Dn}$ by \eqref{eq:pairs-flow-input}, so
$\mathcal N\ge ck^{-1}a^{\gamma Dn}\to\infty$.
We will use punctures separated at a scale $a^{-l}$ much larger than
$R$. The final source-overlap estimate carries a factor $a^{Ql}$,
so we arrange $Ql\leq\gamma Dn$. Put
\[
 \zeta=\gamma/Q,\qquad m=\lfloor\zeta n\rfloor,\qquad l=Dm.
\]
Both $l$ and $k-l$ tend to infinity.

\begin{lemma}[Probability near a puncture]
\label{lem:near-puncture-probability}
For every generation-$n$ node $i$,
\begin{equation}
 \sum_{\substack{j\text{ at generation }n\\d(p_j,p_i)\le a^{-l}}}P_j
 \le Ca^{-\gamma Dm}.
 \label{eq:near-puncture-probability}
\end{equation}
\end{lemma}

\begin{proof}
Fix an ancestor $A$ at generation $n-m$, whose depth is
$k-l=k_0+D(n-m)$. We claim that only boundedly
many descendants $j$ have $d(p_j,p_i)\le a^{-l}$. Split a geodesic word
\[
 v_j^{-1}=w_ju'_j,\qquad |w_j|=l,\qquad |u'_j|=k-l.
\]
The puncture condition and \eqref{eq:inverse-prefix} give
$(p_j|p_i)_e\ge l-C$ and $(v_j^{-1}|p_j)_e\ge k-C$.
Thus $w_j$ lies a bounded distance from the depth-$l$ point on a fixed
ray to $p_i$, leaving boundedly many possibilities.

Reversing and inverting \emph{the same word} gives the geodesic
$v_j=(u'_j)^{-1}w_j^{-1}$. Its depth-$(k-l)$ prefix is therefore
$(u'_j)^{-1}$. Since
\[
 d(x_j,x_A)\le2(B_0+1)r_{k-l},
\]
thinness and \Cref{lem:prefix} put this prefix a bounded distance from
the depth-$(k-l)$ point on a fixed ray to $x_A$. Hence $u'_j$ also has
boundedly many possibilities. Their products determine boundedly
many vertices $v_j$, and the first step of
\Cref{lem:puncture-multiplicity} bounds the nodes assigned to each.
We count group elements here, so nonuniqueness of geodesic spellings
causes no additional multiplicity. The source and inverse prefixes
together account for the full word.

Each such descendant has $P_j\le P_Aa^{-\gamma Dm}$ by
\eqref{eq:pairs-flow-input}. Sum first over these descendants and
then over ancestors, whose probabilities sum to one.
\end{proof}

For each $i\in\mathcal I$, the fraction of its class with puncture
distance at most $a^{-l}$ is bounded by
\[
 \frac{Ca^{-\gamma Dm}}{\mathcal NP_*}
 \le Ck a^{-\gamma Dm}=o(1).
\]
This includes the diagonal. Consequently at least $c\mathcal N^2$
unordered pairs satisfy
\begin{equation}
 d_{ij}:=d(p_i,p_j)>a^{-l}.
 \label{eq:far-pair}
\end{equation}
Call them \emph{far pairs}. Uniformly for these pairs,
$R/d_{ij}\le a^{-(k-l)}\to0$.

\subsection{Comparison on a compact common domain}

Let $K\subset U$ contain all centers, with its closed
$\epsilon$-neighborhood still in $U$. The upper puncture bound gives
\begin{equation}
 d(h_i y,p_i)\ge A_hR/2
 \quad\Longrightarrow\quad
 d(x_i,y)\le D_y\le2C/A_h.
 \label{eq:exterior-preimage}
\end{equation}
Choose $A_h$ so that $2C/A_h<\epsilon/2$ and $A_h/2>2C/d_1$, and set
\[
 E_i=\overline B(p_i,A_hR).
\]
Every $h_iC_{N,t}$ reaches $B(p_i,A_hR/2)$: the diameter bound supplies
a source point at distance at least $d_1/2$ from $x_i$, whose image
has puncture distance at most $2CR/d_1$.

The target balls lie a uniform distance $\delta_*>0$ from $p_i$.
Indeed, their preimages lie in the outer source balls, where
$D_y\le(B_0+2)R$, so the lower puncture bound applies for large $n$.
Because the pattern is shared, this holds for every $i\in\mathcal I$.
For large $n$, uniformly for far pairs,
\[
 2A_hR<a^{-l},\qquad A_hR<\delta_*/2.
\]
Thus $E_i,E_j$ are disjoint, each lies compactly in the other's
transformed domain by \eqref{eq:exterior-preimage}, and the pattern
balls avoid both. These thresholds are independent of $N$.

\begin{lemma}[Pair alternative]
\label{lem:pair-alternative}
Every far pair at a sufficiently large generation has an endpoint
$i$ satisfying
\begin{equation}
 \osc_{h_i^{-1}E_j}u\ge\frac{\kappa}{16}\Delta_i.
 \label{eq:pair-alternative}
\end{equation}
\end{lemma}

\begin{proof}
Write $F_i=u\circ h_i^{-1}$ and $F_j=u\circ h_j^{-1}$ on their open
domains. If $\osc_{E_i}F_j\ge\Delta_j/16$, endpoint $j$ qualifies.
Here we use $\kappa\le1$.
Otherwise, one pattern ball $\mathcal B^s$ has its $F_j$ range interval
at distance at least $\Delta_j/16$ from that on $E_i$. Indeed, if
neither did, the two pattern intervals would have distance less than
$3\Delta_j/16$, contradicting their separation by $\Delta_j/4$.

Node $i$ supplies a compact path $\eta\subset\mathcal B^s$ with
endpoint $F_i$-gap at least $\kappa\Delta_i$. Suppose
$\osc_{E_j}F_i<\kappa\Delta_i/4$. Removing this latter range interval
from the endpoint gap leaves a positive-length closed interval
$I\subset(0,1)$ strictly between the endpoint values and a positive
distance from the range on $E_j$.
To obtain positive margins, choose a smaller closed interval inside
one of the remaining open intervals. Its length need not be uniform
over pairs.

\emph{Fix the pair and $I$ before letting $N\to\infty$.}
Put $F_{i,N}=H_N\circ h_i^{-1}$ and $F_{j,N}=H_N\circ h_j^{-1}$.
The fixed range margins and local uniform convergence imply, for
every $t\in I$ and all sufficiently large $N$, that $h_iC_{N,t}$ hits
$\eta$ and avoids $E_j$. It also reaches the inner $i$-puncture ball.
For the first assertion, the endpoint values still strictly straddle
all of $I$, so \Cref{lem:level-continua} applies to the preimage path.
For the second, the range on $E_j$ remains disjoint from $I$.
By \Cref{lem:stop-continuum}, it has a subcontinuum $K_{N,t}$ containing
a hit on $\eta$, meeting the sphere $d(z,p_i)=A_hR$, and avoiding
$B(p_i,A_hR)$. It remains in the level $F_{i,N}=t$ and avoids $E_j$.
In particular,
\[
 K_{N,t}\subset K_{ij}:=
 Z\setminus\bigl(B(p_i,A_hR)\cup B(p_j,A_hR)\bigr)
 \Subset h_iU\cap h_jU.
\]
Compact containment follows from \eqref{eq:exterior-preimage} with
half the radius. The sphere hit belongs to $E_i$; no identification
of that sphere with the boundary of a metric ball is needed.

Uniform convergence on $\eta\cup E_i$ preserves the $F_j$ range gap,
giving $\osc_{K_{N,t}}F_{j,N}\ge\Delta_j/32$ uniformly for $t\in I$.
The finite-cell comparison, \Cref{lem:cross-energy}, now gives
\[
 \frac{\Delta_j}{32}|I|\le
 \sum_{\substack{q\in T_N\\B(q,2r_N)\cap K_{ij}\ne\varnothing}}
 \osc_{\overline B(q,2r_N)}F_{i,N}\,
 \osc_{\overline B(q,2r_N)}F_{j,N}.
\]
For large $N$ these closed cells lie in the common domain. Both squared
oscillation sums tend to zero by \Cref{lem:local-oscillation}, so
Cauchy--Schwarz contradicts the positive left side. Therefore
$\osc_{E_j}F_i\ge\kappa\Delta_i/4$, proving the alternative.
The argument applies separately to every pair above the geometric
threshold; no convergence rate uniform over pairs is required.
\end{proof}

\subsection{Cost balls and total multiplicity}

We now turn the pair alternative into a measure bound and control the
overlap of the sets being charged. The two indices have distinct roles:
with $i$ fixed, the target punctures control the number of partners $j$;
at a source point, separation of the centers controls the possible
indices $i$. Both estimates will apply to the enlarged balls required
by \eqref{eq:measure}.

For an ordered far pair put $z_{ij}=h_i^{-1}p_j$. The relevant source
scales are
\[
 \begin{array}{c|c}
 \text{puncture denominator }D_{z_{ij}}&\asymp R/d_{ij}\\
 \text{radius of the charged neighborhood}&R^2/d_{ij}^2.
 \end{array}
\]
The radius divided by the denominator is $O(R/d_{ij})\to0$. This separation of scales is what
allows a fixed enlargement of the charged neighborhood without losing
control of its image.

If $y\in h_i^{-1}E_j$, then
$d(h_i y,p_i)\ge d_{ij}-A_hR\ge d_{ij}/2$, whence
$D_y\le CR/d_{ij}$. These points approach $x_i$ uniformly, so the
two-sided puncture bound applies and gives
\begin{equation}
 D_{z_{ij}}\asymp R/d_{ij}.
 \label{eq:cost-denominator}
\end{equation}
The lower expansion bound and $d(h_i y,p_j)\le A_hR$ then give
$d(y,z_{ij})\le CR^2/d_{ij}^2$. Fixing a sufficiently large $C_2$, put
$s_{ij}=C_2R^2/d_{ij}^2$; thus
\begin{equation}
 h_i^{-1}E_j\subseteq\overline B(z_{ij},s_{ij}).
 \label{eq:cost-containment}
\end{equation}
Choose the enlargement constant $C_\mu$ in \eqref{eq:measure}
uniformly on a compact neighborhood of the centers and define
\[
 \mathcal D_{ij}=\overline B(z_{ij},C_\mu s_{ij}).
\]
Their centers stay in this neighborhood and their radii tend uniformly
to zero. If $i$ satisfies \Cref{lem:pair-alternative}, \eqref{eq:measure} gives
\begin{equation}
 \mu(\mathcal D_{ij})\ge c\kappa^2\Delta_i^2.
 \label{eq:cost-lower-bound}
\end{equation}

The enlarged balls still have small images. Indeed, $D_y$ is
$1$-Lipschitz and
\[
 \frac{C_\mu s_{ij}}{D_{z_{ij}}}
 \le C\frac R{d_{ij}}\le Ca^{-(k-l)}\longrightarrow0.
\]
Thus $D_y\asymp D_{z_{ij}}$ on $\mathcal D_{ij}$, and the upper
expansion bound yields
\begin{equation}
 h_i\mathcal D_{ij}\subseteq\overline B(p_j,C_3R).
 \label{eq:cost-image}
\end{equation}
Figure~\ref{fig:pair-cost} shows the charged neighborhood and the image
control that survives the fixed enlargement.

\begin{figure}[htbp]
\centering
\begin{tikzpicture}[x=1cm,y=1cm,font=\small,
  line cap=round,line join=round,>=Stealth]
  \definecolor{costblue}{RGB}{54,101,125}
  \definecolor{costochre}{RGB}{164,111,48}
  \node[font=\bfseries] at (2.75,4.95) {Source};
  \node[font=\bfseries] at (11.40,4.95) {Target};

  \fill[costblue!5] (3.25,2.80) circle (1.10);
  \draw[costblue,thick] (3.25,2.80) circle (1.10);
  \draw[costblue,dashed,thick] (3.25,2.80) circle (.52);
  \node[costblue,anchor=south] at (3.25,4.00) {$\mathcal D_{ij}$};
  \fill[black!70] (.55,2.80) circle (1.8pt);
  \node[anchor=north] at (.55,2.54) {$x_i$};
  \draw[black!40,thin] (3.25,2.73)--(3.25,1.47);
  \fill[costblue] (3.25,2.80) circle (1.8pt);
  \node[anchor=north] at (3.25,1.37) {$z_{ij}$};
  \node[align=center] at (2.75,.57)
    {$h_i^{-1}E_j\subseteq\overline B(z_{ij},s_{ij})$};
  \node[align=center] at (2.75,.05)
    {$s_{ij}\asymp R^2/d_{ij}^{\,2}$};

  \draw[black!85,thick,->] (5.55,2.80)--(7.25,2.80)
    node[midway,above=5pt] {$h_i$};

  \fill[costochre!7] (9.00,2.80) circle (.50);
  \fill[costochre!7] (13.20,2.80) circle (.50);
  \draw[costochre,thick] (9.00,2.80) circle (.50);
  \draw[costochre,thick] (13.20,2.80) circle (.50);
  \draw[costblue,dashed,thick] (13.20,2.80) circle (1.10);
  \node[costochre,anchor=south] at (9.00,3.57) {$E_i$};
  \node[costochre,anchor=west] at (14.50,3.65) {$E_j$};
  \draw[costochre!75,thin] (14.38,3.57)--(13.51,3.19);
  \draw[black!40,thin] (9.00,2.73)--(9.00,1.47);
  \draw[black!40,thin] (13.20,2.73)--(13.20,1.47);
  \fill[costochre] (9.00,2.80) circle (1.8pt);
  \fill[costochre] (13.20,2.80) circle (1.8pt);
  \node[anchor=north] at (9.00,1.37) {$p_i$};
  \node[anchor=north] at (13.20,1.37) {$p_j$};
  \draw[black!85,<->] (9.00,.76)--(13.20,.76)
    node[midway,below=3pt] {$d_{ij}\gg R$};
  \node[costblue,align=center] at (11.10,-.30)
    {$h_i\mathcal D_{ij}\subseteq\overline B(p_j,C_3R)$};
\end{tikzpicture}
\caption{The measure cost of a far pair. The dashed source ball contains
$h_i^{-1}E_j$ and has radius $s_{ij}\asymp R^2/d_{ij}^2$.
Its fixed enlargement $\mathcal D_{ij}$, drawn solid, still maps into
the dashed target ball of radius $C_3R$. Here $p_i=h_i z_0$, $z_{ij}=h_i^{-1}p_j$, and
$E_i,E_j$ are the solid puncture balls of radius $A_hR$.
The circles are schematic; no connectedness or boundary regularity of
metric balls is assumed.}
\label{fig:pair-cost}
\end{figure}

For fixed $i$, \Cref{lem:puncture-multiplicity} therefore bounds the
multiplicity of $\mathcal D_{ij}$ as $j$ varies. Also
\begin{equation}
 \mathcal D_{ij}\subseteq\overline B(x_i,C_4Ra^l),
 \label{eq:cost-source}
\end{equation}
because $d(x_i,z_{ij})\le CR/d_{ij}$ and the cost radius is
$o(R/d_{ij})$. The $HR$-separation of the $x_i$ and
\Cref{lem:packing} bound the multiplicity of these larger balls by
$Ca^{Ql}$.

\begin{proof}[Conclusion of the proof of \Cref{prop:no-subexponential-records}]
Assign each unordered far pair to a qualifying endpoint. At any point,
\eqref{eq:cost-source} allows only $Ca^{Ql}$ paying indices $i$, and
\eqref{eq:cost-image} allows boundedly many partners for each. Integrating
this multiplicity bound gives
\begin{equation}
 \sum_{\{i,j\}\text{ far, assigned to }i}\mu(\mathcal D_{ij})
 \le Ca^{Ql}\mu(Z).
 \label{eq:cost-total-upper}
\end{equation}
This counts the enlarged balls directly and uses no invariance of $\mu$.
The assignment may be arbitrary when both endpoints qualify: all
lower bounds hold for either endpoint, and the multiplicity estimate
already counts every possible ordered contribution.

Conversely, every paying endpoint satisfies
\[
 \Delta_i^2\ge\Delta_0^2a^{2\gamma Dn}P_i^2
 \ge\Delta_0^2a^{2\gamma Dn}P_*^2.
\]
There are at least $c\mathcal N^2$ far pairs, so
\eqref{eq:cost-lower-bound} and \eqref{eq:pairs-class-size} give total cost
\begin{equation}
 c\mathcal N^2\kappa^2\Delta_0^2a^{2\gamma Dn}P_*^2
 \ge c'k^{-2}a^{2\gamma Dn}.
 \label{eq:cost-total-lower}
\end{equation}
Since $\kappa,\Delta_0>0$ are fixed and
$Ql=QD\lfloor\gamma n/Q\rfloor\le\gamma Dn$, comparison with
\eqref{eq:cost-total-upper} yields
\[
 a^{\gamma Dn}\le Ck^2.
\]
This contradicts $k=k_0+Dn$, excluding the unbounded subexponential case.
\end{proof}


\section{Completion of the proof}
\label{sec:completion}

\begin{proof}[Proof of \Cref{thm:main}]
If $M(n)$ were unbounded, the record selection of \Cref{sec:scalar}
would give either exponential records, excluded by
\Cref{prop:no-exponential-records}, or subexponential growth,
excluded by \Cref{prop:no-subexponential-records}. Thus $\sup_n M(n)<\infty$.
\Cref{thm:uniformization} uniformizes the boundary, and
\Cref{prop:action} supplies the proper cocompact action with finite kernel.
\end{proof}

\section{Consequences}
\label{sec:consequences}

\begin{corollary}
\label{cor:torsion-free}
If $G$ is a torsion-free hyperbolic group and $\partial G$ is homeomorphic to $\Sphere$,
then $G$ is the fundamental group of a closed hyperbolic three-manifold.
\end{corollary}
\begin{proof}
Point stabilizers in a proper action are finite, hence trivial when
$G$ is torsion-free. Thus the action in \Cref{thm:main} is free, and
$\HH^3/G$ is a closed hyperbolic three-manifold with fundamental group
$G$. The manifold may be nonorientable.
\end{proof}

\begin{corollary}[Virtual structure in the torsion-free case]
\label{cor:virtual-structure}
Let $G$ be a torsion-free hyperbolic group with $\partial G$ homeomorphic
to $\Sphere$. Some finite-index subgroup of $G$ is isomorphic to
$\pi_1(\Sigma)\rtimes\ZZ$ for a closed connected orientable surface
$\Sigma$. In particular, $G$ has positive virtual first Betti number.
Moreover, $G$ is virtually compact special and residually finite.
\end{corollary}
\begin{proof}
Let $M^+$ be the orientation cover of the closed hyperbolic manifold in
\Cref{cor:torsion-free}, using the manifold itself when it is orientable,
and let $G^+=\pi_1(M^+)$. Thus $G^+$ is normal of index at most two in
$G$. Agol's virtual-fibering theorem
\cite[Theorem 9.2]{Agol} gives a connected finite cover $N$ of $M^+$
fibering over the circle with closed connected orientable fiber $\Sigma$.
The fibration sequence gives
$\pi_1(N)\cong\pi_1(\Sigma)\rtimes\ZZ$, with action induced by the
monodromy, and its surjection onto $\ZZ$ gives $b_1(\pi_1(N))\ge1$.

Bergeron--Wise's cubulation theorem \cite[Theorem 5.3]{BergeronWise}
gives a proper cocompact action of $G^+$ on a CAT(0) cube complex, free
because $G^+$ is torsion-free. Agol's virtual-special theorem
\cite[Theorem 1.1]{Agol} then gives a finite-index subgroup of $G^+$
with compact special quotient, so $G$ is virtually compact special.
Agol's separability corollary \cite[Corollary 1.3]{Agol}, applied to the
trivial subgroup, also makes $G^+$ residually finite. The finite quotient
$G/G^+$ and intersections of the finitely many $G$-conjugates of
finite-index normal subgroups of $G^+$ then separate every nontrivial
element of $G$.
\end{proof}

\begin{thebibliography}{99}

\bibitem{Agol}
I.~Agol, The virtual Haken conjecture, with an appendix by
I.~Agol, D.~Groves, and J.~Manning,
\emph{Doc. Math.} \textbf{18} (2013), 1045--1087.
\href{https://doi.org/10.4171/DM/421}{doi:10.4171/DM/421}.

\bibitem{AhlforsBers}
L.~Ahlfors and L.~Bers, Riemann's mapping theorem for variable metrics,
\emph{Ann. of Math.} (2) \textbf{72} (1960), 385--404.
\href{https://doi.org/10.2307/1970141}{doi:10.2307/1970141}.

\bibitem{AlbinPoggiCorradini}
N.~Albin and P.~Poggi-Corradini,
Minimal subfamilies and the probabilistic interpretation for modulus on graphs,
\emph{J. Anal.} \textbf{24} (2016), 183--208.
\href{https://doi.org/10.1007/s41478-016-0002-9}{doi:10.1007/s41478-016-0002-9}.

\bibitem{AIM}
K.~Astala, T.~Iwaniec, and G.~Martin,
\emph{Elliptic Partial Differential Equations and Quasiconformal
Mappings in the Plane}, Princeton Mathematical Series 48,
Princeton University Press, 2009.

\bibitem{BergeronWise}
N.~Bergeron and D.~T. Wise, A boundary criterion for cubulation,
\emph{Amer. J. Math.} \textbf{134} (2012), 843--859.
\href{https://doi.org/10.1353/ajm.2012.0020}{doi:10.1353/ajm.2012.0020}.

\bibitem{BK}
M.~Bonk and B.~Kleiner,
Quasisymmetric parametrizations of two-dimensional metric spheres,
\emph{Invent. Math.} \textbf{150} (2002), 127--183.
\href{https://doi.org/10.1007/s00222-002-0233-z}{doi:10.1007/s00222-002-0233-z}.

\bibitem{BKConformalDimension}
M.~Bonk and B.~Kleiner,
Conformal dimension and Gromov hyperbolic groups with 2-sphere boundary,
\emph{Geom. Topol.} \textbf{9} (2005), 219--246.
\href{https://doi.org/10.2140/gt.2005.9.219}{doi:10.2140/gt.2005.9.219}.

\bibitem{BourdonKleiner}
M.~Bourdon and B.~Kleiner,
Combinatorial modulus, the combinatorial Loewner property, and Coxeter groups,
\emph{Groups Geom. Dyn.} \textbf{7} (2013), 39--107.
\href{https://doi.org/10.4171/GGD/177}{doi:10.4171/GGD/177}.

\bibitem{BH}
M.~R. Bridson and A.~Haefliger,
\emph{Metric Spaces of Non-Positive Curvature},
Grundlehren der mathematischen Wissenschaften 319, Springer, 1999.
\href{https://doi.org/10.1007/978-3-662-12494-9}{doi:10.1007/978-3-662-12494-9}.

\bibitem{Cannon}
J.~W. Cannon, The combinatorial Riemann mapping theorem,
\emph{Acta Math.} \textbf{173} (1994), 155--234.
\href{https://doi.org/10.1007/BF02398434}{doi:10.1007/BF02398434}.

\bibitem{CFPSquaring}
J.~W. Cannon, W.~J. Floyd, and W.~R. Parry,
Squaring rectangles: the finite Riemann mapping theorem,
in \emph{The Mathematical Legacy of Wilhelm Magnus: Groups, Geometry
and Special Functions}, Contemporary Mathematics 169,
American Mathematical Society, 1994, 133--212.
\href{https://personal.math.vt.edu/floyd/research/papers/mageps.pdf}{Author text}.

\bibitem{CFP}
J.~W. Cannon, W.~J. Floyd, and W.~R. Parry,
Sufficiently rich families of planar rings,
\emph{Ann. Acad. Sci. Fenn. Math.} \textbf{24} (1999), 265--304.
\href{https://www.acadsci.fi/mathematica/Vol24/cannon.pdf}{Journal text}.

\bibitem{CFPFiniteSubdivision}
J.~W. Cannon, W.~J. Floyd, and W.~R. Parry,
Finite subdivision rules,
\emph{Conform. Geom. Dyn.} \textbf{5} (2001), 153--196.
\href{https://doi.org/10.1090/S1088-4173-01-00055-8}{doi:10.1090/S1088-4173-01-00055-8}.

\bibitem{CannonSwenson}
J.~W. Cannon and E.~L. Swenson,
Recognizing constant curvature discrete groups in dimension 3,
\emph{Trans. Amer. Math. Soc.} \textbf{350} (1998), 809--849.
\href{https://doi.org/10.1090/S0002-9947-98-02107-2}{doi:10.1090/S0002-9947-98-02107-2}.

\bibitem{Drilling}
D.~Groves, P.~Ha\"issinsky, J.~F. Manning, D.~Osajda,
A.~Sisto, and G.~S. Walsh, Drilling hyperbolic groups,
preprint, 2026.
\href{https://arxiv.org/abs/2406.14667v2}{arXiv:2406.14667v2}.

\bibitem{Haissinsky}
P.~Ha\"issinsky, Hyperbolic groups with planar boundaries,
\emph{Invent. Math.} \textbf{201} (2015), 239--307.
\href{https://doi.org/10.1007/s00222-014-0552-x}{doi:10.1007/s00222-014-0552-x}.

\bibitem{HK}
J.~Heinonen and P.~Koskela, Definitions of quasiconformality,
\emph{Invent. Math.} \textbf{120} (1995), 61--79.
\href{https://doi.org/10.1007/BF01241122}{doi:10.1007/BF01241122}.

\bibitem{KahnMarkovic}
J.~Kahn and V.~Markovi\'c,
Immersing almost geodesic surfaces in a closed hyperbolic three manifold,
\emph{Ann. of Math.} (2) \textbf{175} (2012), 1127--1190.
\href{https://doi.org/10.4007/annals.2012.175.3.4}{doi:10.4007/annals.2012.175.3.4}.

\bibitem{Markovic}
V.~Markovi\'c, Criterion for Cannon's conjecture,
\emph{Geom. Funct. Anal.} \textbf{23} (2013), 1035--1061.
\href{https://doi.org/10.1007/s00039-013-0228-5}{doi:10.1007/s00039-013-0228-5}.

\bibitem{Pansu1989}
P.~Pansu, Dimension conforme et sph\`ere \`a l'infini des vari\'et\'es
\`a courbure n\'egative,
\emph{Ann. Acad. Sci. Fenn. Ser. A I Math.} \textbf{14} (1989), 177--212.
\href{https://doi.org/10.5186/aasfm.1989.1424}{doi:10.5186/aasfm.1989.1424}.

\bibitem{RajalaUniformization}
K.~Rajala, Uniformization of two-dimensional metric surfaces,
\emph{Invent. Math.} \textbf{207} (2017), 1301--1375.
\href{https://doi.org/10.1007/s00222-016-0686-0}{doi:10.1007/s00222-016-0686-0}.

\bibitem{Schramm}
O.~Schramm, Square tilings with prescribed combinatorics,
\emph{Israel J. Math.} \textbf{84} (1993), 97--118.
\href{https://doi.org/10.1007/BF02761693}{doi:10.1007/BF02761693}.

\bibitem{Sullivan}
D.~Sullivan, On the ergodic theory at infinity of an arbitrary discrete
group of hyperbolic motions, in \emph{Riemann Surfaces and Related Topics:
Proceedings of the 1978 Stony Brook Conference}, I.~Kra and B.~Maskit, eds.,
Annals of Mathematics Studies 97, Princeton University Press, 1981, 465--496.
\href{https://doi.org/10.1515/9781400881550-035}{doi:10.1515/9781400881550-035}.

\bibitem{Tukia}
P.~Tukia, On two-dimensional quasiconformal groups,
\emph{Ann. Acad. Sci. Fenn. Ser. A I Math.} \textbf{5} (1980), 73--78.
\href{https://doi.org/10.5186/aasfm.1980.0530}{doi:10.5186/aasfm.1980.0530}.

\end{thebibliography}

\end{document}
